% ============================================================
%  Georg Brandes, Samlede Skrifter. Trettende Bind (Supplementbind).
%  Kjøbenhavn: Gyldendalske Boghandels Forlag (F. Hegel & Søn).
%  Græbes Bogtrykkeri.  1903.
%  Indledning (1903) and "Striden om Tro og Viden" (1866--67).
%  TRANSCRIPTION (Danish), printed pp. 1--42.
%
%  One of three transcription directories cut from this volume:
%      striden/               pp. 1--42    Indledning; Striden om Tro og Viden
%      nielsen-anmeldelser/   pp. 85--109  five pieces on Nielsen, 1867
%      brochner-bidrag/       pp. 115--123 review of Brøchner's Bidrag (1869)
%  NOT transcribed: pp. 43--84 (Dualismen i vor nyeste Filosofi), 110--114,
%  124 ff.  See the note on pp. 43--84 below.
%
%  ---- DATE -------------------------------------------------------------
%  The title page (PDF 7) reads 1903.  The Indledning (p. 3) is signed
%  "April 1903." / "G. B."  (catalog.yaml's 1899 is the series date, not
%  this volume's.)  This is Brandes's 1903 revision of pieces first printed
%  1865--69, in the orthography of 1903 (Filosofi, Teologi, eksakt, hv-,
%  single vowels), transcribed as printed.  Nothing has been harmonised
%  with the first printings or with the 1866 Dualismen.
%
%  ---- SOURCE -----------------------------------------------------------
%  bibliotek/Brandes, Georg/1903-samlede-skrifter-13.pdf
%      Google Books (University of Chicago copy), 552 PDF pp., 600 dpi
%      bitonal JBIG2, embedded text layer (witness 1).
%      sha256 9579f1fd8672884bcc7e149b43f5c8c640dcaaefa1aceca01d5f89cbed016704
%  Witness 2 (never copy-text): Project Runeberg's OCR text of the same
%      edition (gbsamskr/13), bibliotek/Brandes, Georg/gbsamskr-13.txt,
%      cut into printed pages by alignment against witness 1 (per-page word
%      agreement between the two witnesses, pp. 1--42: mean 0.966, min 0.935 over 39 text pages).
%      It is a different OCR engine's reading; whether Runeberg photographed
%      a different physical copy is not known.
%
%  ---- PAGE MAP ---------------------------------------------------------
%      printed p. = PDF p. - 12, uniform over printed 1--136.
%  Verified 2026-09-22 AT THE IMAGE: folio read off the rendered page at
%  printed 9, 29, 38, 42, 44, 86, 97, 104, 109, 116, 123 (all match; pp. 2
%  and 17 print no folio and were placed by content).  Independently read
%  by H. Halvorson at 8, 9, 42, 43, 54, 85, 93, 123, 136.
%
%  ---- LAYOUT OF THE VOLUME ---------------------------------------------
%  Every piece begins on a fresh page with a sunk head: title in large
%  italic, the year of first publication in roman in parentheses, a short
%  rule; the opening page carries no folio and no running head.  Every piece
%  ends with a short centred rule (\piecerule).  Running heads carry a short
%  title of the piece (a third reading, recorded at each piece's head where
%  it differs); they are not transcribed.
%
%  ---- CONVENTIONS ------------------------------------------------------
%  Diplomatic.  Orthography, punctuation and capitalisation as printed.
%  Emphasis: ITALIC is the volume's only emphasis device -- no letterspacing
%    and no bold were found in pp. 1--123 -- and every italic span is set
%    \emph{}, whether it marks emphasis, a title or a foreign word; no
%    attempt is made to distinguish the functions.
%  Quotation marks: the volume quotes with SINGLE guillemets pointing
%    outward, ‹...› (U+2039/U+203A), for quotations, titles and terms alike;
%    double «...» occur only in a head (p. 102).  Transcribed as printed.
%  Dashes: spaced em-dash, set ---.  ɔ: as printed.
%  Footnotes: *), set with \fnmark{*)}{...} at the mark.
%  A word broken across a printed line is joined (Cau\-%); a hyphen printed
%    at a line end as part of the word is kept.
%  Printer's errors are transcribed AS PRINTED and logged in a % comment at
%    the site; nothing is silently corrected.
%  Punctuation is as the Chicago copy prints it.  The 600 dpi bitonal image
%    sometimes loses the tail of a faint comma; where it shows a point but
%    witness 1 (Google's OCR, made from its capture BEFORE compression to
%    bitonal JBIG2) reads a comma, the comma is transcribed and commented;
%    where the image and witness 1 both show a point, the point is
%    transcribed and logged as a printer's error, whatever the syntax wants.
%  Letters the bitonal scan may not resolve (ø/o, æ/œ, worn e/c): where the
%    printed form would not be a word, the word is read unless the glyph
%    positively shows otherwise (TRANSCRIPTION-PLAYBOOK §2 rule 3); each case
%    is commented.  In this fount the italic æ has a closed o-like bowl and no
%    separate a-stem, so it looks like œ (see p. 27 `duæ', p. 10 Fædrelandet).
%  A hyphen at a line end that could be either a compound hyphen or a word
%    break is resolved by the volume's own usage MID-LINE: printed only solid
%    -> joined; only hyphenated -> kept; both or never -> the printed hyphen
%    is kept.  Each case is commented with its evidence.
%  A word divided across a PAGE turn is joined with `%' before the page
%    marker (e.g. Spørgs%  / \opage{122}maal); the divisional hyphen is not set.
%
%  ---- CONTENTS vs HEADS ------------------------------------------------
%  The Indhold (PDF 9) gives short titles; the heads over the pieces name
%  the book under review and the year.  Both readings are kept: the Indhold
%  reading is the first argument of \piecehead (table of contents, running
%  head), the head as printed is the second.  Recorded at each piece.
%
%  ---- READINGS OF THE HEADS ----------------------------------------------
%  Indhold (PDF 9)                       Head as printed over the piece
%  Indledning 1                          INDLEDNING (roman caps); signed
%                                          `April 1903.' / `G. B.' (p. 3)
%  Striden om Tro og Viden 5 (italic)    STRIDEN OM TRO OG VIDEN, section
%                                          title (p. 5); pp. 4, 6 blank
%  Samvittighed og Videnskab 7           A. C. Larsen: Samvittighed og
%                                          Videnskab (1866) -- p. 10 dated
%                                          `I December 1865.'
%  Svar til Hr. Schmidt 11               Svar til Hr. Schmidt (1866)
%  Filosofi og Teologi 13                I Anledning af Hr. Harald Høffdings
%                                          / Stridsskrift: Filosofi og Teologi (1866)
%  En sidste Replik 17                   En sidste Replik (1866)
%  Om Selvansvarlighed 22                Om Selvansvarlighed (1866)
%  Nationalitetsideens Gyldighed 26      Claudius Wilckens: Nationalitetsideens
%                                          Gyldighed (1867)
%  Indledning til den rationelle Etik 28 P. S. V. Heegaard: Indledning til den
%                                          rationelle Etik (1867)
%  Gensvar til Dr. Heegaard 34           Gensvar til Hr. Dr. P. S. V. Heegaard (1867)
%  Grundideernes Logik 40                R. Nielsen: Grundideernes Logik
%                                          -- NO year printed (300/600 dpi)
%
%  ---- PRINTER'S ERRORS AND DOUBTFUL READINGS (each commented at its site) --
%  p. 2 `1889---90' printed with an unspaced dash; a stray dot before
%  `misforstaas' (a speck, not set).  p. 3 `dèr' with grave.  p. 11 broken e
%  in `er'.  p. 13 opening guillemet turned: ›Samvittighed og Videnskab›.
%  p. 15 `Noget,' and p. 20 `den,': faint marks, witness 1 reads commas.
%  p. 16 `Heffdings' (wrong-fount e); `bebuder.' point for comma.  p. 17
%  `at kan kom' (kan for han).  p. 18 note: ‹at slaa ... never closed.
%  p. 20 `Menne-/sker' printed without its line-end hyphen; note: `Løsning',
%  slash not visible (sense reading).  p. 23 closing guillemet turned:
%  vedligeholdes‹.  p. 24 no stop after `Forvirring'.  p. 26, 30 broken
%  type.  p. 27 `duæ' (see Conventions).  p. 30 no stop after `Aandsmagt'.
%  p. 31 `etisko'.  p. 37 `indrømme. at'.  p. 38 `Saa frejdig kan ellers
%  er' (kan for han).  p. 41 `‹Paradoxet›' with x (elsewhere k), as printed.
%  Guillemet balance ‹ minus › = +1, and every unit of it is one of the
%  logged defects above (p. 13 -2, p. 18 +1, p. 23 +2).
%
%  ---- pp. 43--84, NOT TRANSCRIBED --------------------------------------
%  pp. 43--84 reprint "Dualismen i vor nyeste Filosofi" (1866).  The 1866
%  first edition is already in this collection (texts/brandes/dualismen),
%  fully transcribed and collated against two scans.  The 1903 text is
%  Brandes's own revision, in modernised orthography and with altered
%  chapter titles (Indhold: I. Indledning 43, Den absolute Uensartethed 54,
%  Det Antirationelle 67, Dualismens Følger 71), so it is a VARIANT WITNESS
%  to the 1866 text and deserves a collation of its own, not a second
%  transcription.
%
%  ---- HOW THIS TEXT'S ACCURACY WAS ESTABLISHED ------------------------
%  Word-accuracy: every batch diffed against TWO machine witnesses before
%  splicing (ocrdiff.py; TRANSCRIPTION-PLAYBOOK §3 step 4) -- witness 1 the
%  embedded layer, witness 2 Runeberg's OCR -- and then checked with
%  consensus.py, which compares case-sensitively and with punctuation where
%  the two witnesses agree with each other (ocrdiff lower-cases and ignores
%  punctuation).  4 batches; witness 1: 10 candidates, witness 2: 32,
%  in both: 2, consensus: 6; 0 transcription errors surfaced; all were
%  the witnesses' faults; 0 unresolved.
%  Independent collation: a seeded random sample of 8 of the 39 text pages
%  (seed 20260922: pp. 1, 3, 8, 10, 27, 36, 37, 40) was read line by line at
%  the image by a reader who had not transcribed them -- words, capitals,
%  punctuation, italics, paragraphing, page turns.  1 discrepancy found
%  (p. 27 `duœ' for the printed `duæ', a deliberate misreading of this
%  fount's italic æ) and corrected: 0.13 errors/page in the sample before
%  correction.  The same reader re-adjudicated every doubtful call the batch
%  agents reported and overturned one (p. 16 `bebuder,' -> `bebuder.', see
%  the site); four italic calls spot-checked, all confirmed.
%  Unmeasured beyond that sample: a sample of 8 pages bounds the rate only
%  loosely (1 error in 8 pages; 95% upper bound ~0.6/page).
% ============================================================

\documentclass[12pt,a4paper,openany]{book}
\usepackage[utf8]{inputenc}
\usepackage[T1]{fontenc}
\usepackage{libertinus}
\usepackage{libertinust1math}
\usepackage{graphicx}    % for \reflectbox (the old Danish ɔ: id-est mark)
\DeclareUnicodeCharacter{0254}{\reflectbox{c}}  % ɔ (U+0254) = reversed c
\usepackage[danish]{babel}
\usepackage[protrusion=true,expansion=true]{microtype}
\usepackage{setspace}
\usepackage[margin=1.2in]{geometry}
\usepackage{fancyhdr}
\setlength{\headheight}{28pt}

% Footnotes as the 1903 printing sets them: *), not 1, 2, 3.  Set the mark
% explicitly at each site with \fnmark{*)}{...} (or \fnmark{**)}{...} should a
% page carry a second note); the default is restored afterwards.
\renewcommand{\thefootnote}{*)}
\newcommand{\fnmark}[2]{\renewcommand{\thefootnote}{#1}\footnote{#2}%
  \renewcommand{\thefootnote}{*)}}

% A piece's head.  #1 = the title as the volume's INDHOLD prints it (used for
% the table of contents and the running head); #2 = the head block as printed
% over the piece, with its own \opage.  The two readings differ, and both are
% kept (TRANSCRIPTION-PLAYBOOK §6).  The short rule under every head is
% printed, and is set by the macro.
\newcommand{\piecehead}[2]{\clearpage\phantomsection
  \addcontentsline{toc}{section}{#1}\markboth{#1}{#1}%
  \begin{center}\large #2\\[6pt]\rule{2cm}{0.4pt}\end{center}\par\medskip}
% The short centred rule printed at the end of every piece.
\newcommand{\piecerule}{\par\begin{center}\rule{1.5cm}{0.4pt}\end{center}\par}

\usepackage{hyperref}
\hypersetup{pdftitle={Brandes, Samlede Skrifter XIII: Striden om Tro og Viden}, pdfauthor={Georg Brandes}, hidelinks}

\pagestyle{fancy}
\fancyhf{}
\fancyhead[LE,RO]{\thepage}
\fancyhead[RE]{\textit{Samlede Skrifter XIII (1903)}}
\fancyhead[LO]{\textit{\leftmark}}
\setstretch{1.3}

\title{\textbf{Striden om Tro og Viden}\\[8pt]\large Georg Brandes, \textit{Samlede Skrifter}, Trettende Bind (Supplementbind),\\ pp.~1--42}
\author{}
\date{Kjøbenhavn: Gyldendalske Boghandels Forlag (F.\ Hegel \& Søn), 1903}

\usepackage{marginnote}
\newcommand{\opage}[1]{\marginnote{\footnotesize\textit{[#1]}}}
\newcommand{\apage}[1]{\marginnote{\footnotesize\textit{[#1]}}}
\begin{document}
\maketitle
\thispagestyle{empty}
\tableofcontents

\mainmatter

% ============================================================
%  Indledning (printed pp. 1--3)
%  Indhold (PDF 9) reads: `Indledning 1'
%  Head as printed: `INDLEDNING' (roman capitals, no year), short rule.
%  First word opens with an enlarged initial `I' (not a drop cap); set plain.
%  Signed `April 1903.' / `G. B.' (right), then a broken short rule.
%  p. 1 foot carries the signature line `G. Brandes: Samlede Skrifter. XIII.'
%  and p. 3 the signature `1*' -- not transcribed.
% ============================================================
% --- p. 1 ---
\piecehead{Indledning}{\opage{1}INDLEDNING}
I Forordet til mine \emph{Samlede Skrifter} blev det bemærket, at
Udgaven ikke vilde indbefatte Alt hvad jeg i mit Liv havde
offenliggjort. Den blev endda langt omfangsrigere end det efter
den oprindelige Beregning havde tegnet til, og indbefattede
bl. A. foruden tidligere samlede Ting 1200 Sider, der endnu
ikke var udgivne i Bogform.

Ikke medtaget blev Alt, hvis Form var polemisk, Alt, hvad
der ikke var umiddelbart tilgængeligt for Læseren, og en Mang\-%
foldighed af kortere Artikler, hvis Optagelse vilde have bragt
Værket til at svulme altfor stærkt.

En Forfatter er imidlertid uheldigt stillet, hvis han ikke
selv samler og tager Vare paa hvad der --- uanset dets større
eller ringere Værd --- bærer Præget af hans Aandsform. Ti,
som Erfaringen viser, hvad han selv lader slippe ud af sine
Hænder, det kaster andre Hænder sig over. Er en Skribent
naaet til at gøre sit Navn tilstrækkeligt bekendt, saa roder
fremmede Fingre i hans literære Levned og fremstiller dette
saaledes, som formentlig Overlegenhed finder for godt. Og denne
Overlegenhed har des friere Spil, jo mere han har overladt til
dens Vilkaarlighed.

Denne Synsmaade har bestemt mig til at udgive et Supplement
til mine Samlede Skrifter, og jeg aabner dette første Supplement%
% --- p. 2 ---
\opage{2}bind med mine tidligste polemiske Artikler, paa hvilket et Ud\-%
valg af andre vil følge, blandede med Arbejder, der kun vil
fremstille og belære. Med Hensyn til disse ældste Artikler og
de Par Spydigheder, de indeholder mod enkelte udmærkede
Forfattere som Professorerne Høffding og Wilckens, hvis Vel\-%
vilje er mig en Ære, jeg ikke kan mistænkes for at ville for\-%
spilde, og hvem jeg derfor mindst af Alt vilde krænke, maa det
være mig tilladt at minde om det smukke Brev fra J. L. Hei\-%
berg til Carsten Hauch (28 Januar 1841), i hvilket denne efter
med Hjertelighed at have modtaget Hauchs udstrakte Haand
og forsikret ham, at en gammel Fejde imellem dem ikke har
efterladt Gnist af Uvilje i hans Sind, siger ham, at et Optryk
af nogle polemiske Artikler fra hin Tid ‹ingenlunde kan eller
bør forstyrre det gode Forhold›, som til hans Glæde nu bestaar
dem imellem: ‹Jeg véd vel, at denne Fornyelse af hvad man i
% a stray point is printed on the baseline between `lettelig' and
% `misforstaas' (a displaced or foul full stop, or dirt); not transcribed.
én Henseende kunde ønske glemt, lettelig misforstaas›. Men,
siger han ‹det eksisterer nu engang, og i Literaturen saa vel
som i den politiske Historie kan Fakta ikke alene ikke for\-%
andres, men det er endog Pligt at samle og opbevare dem.›
Ganske saa vidt som J. L. Heiberg gaar ikke jeg. Mangfoldige
polemiske Enkeltheder overgives bedst til Glemselen, og det kunde
f. Eks. ikke falde mig ind at genoptrykke mine polemiske Ud\-%
% `1889---90': the range is printed with an unspaced em-dash.
fald fra Nietzsche-Fejden 1889---90 efter at Høffding har udeladt
sine af Samlingen \emph{Mindre Arbejder}. Men til nogle af de literære
Fejder, hvori jeg er bleven indviklet, vil man maaske i Frem\-%
tiden vende tilbage. Meget, som nu synes kun at have historisk
Interesse, kan, naar beslægtede Foreteelser atter opstaar, paany
blive aktuelt.

Overfor en kritisk Forfatter deler Læseverdenen sig i to
Grupper, en større, der læser ham, ifald den interesserer sig for
de Stoffer, han behandler, og en mindre, der interesserer sig
for Forfatteren selv. I de \emph{Samlede Skrifter} var, som det i For\-%
ordet udtrykkelig blev sagt, hele Eftertrykket lagt paa det Stoflige.
% --- p. 3 ---
% `dèr' is printed with a grave accent (witness 2 reads `dér').
% `Æmnernes', `Æmnet', `Æmne' as printed (witness 1 reads `Emn-').
\opage{3}Hensynet til det store Publikums Trang til Belæring var det
ene afgørende dèr. Hvad der ved Æmnernes Art hørte sammen,
blev aldrig splittet ad, var det end frembragt paa vidt fra hin\-%
anden liggende Tider. Og Intet var medtaget, som kun ved
Behandlingsmaaden kunde gøre Krav paa nogen Opmærksomhed.

I Supplementet skulde gerne ogsaa Et og Andet finde Plads,
som kun kan fængsle dem, der forud nærer Interesse for den
Skrivende og som læser ham, selv om Æmnet han behandler,
er dem ligegyldigt. Den kritiske Fremstilling af en Personlighed
er jo en selvstændig Kunst og kan vurderes omtrent som Por\-%
trætmaleriet uden Hensyn til om Læser eller Beskuer kender
Genstanden eller ej. Fægtekunsten i aandelige Tvistemaal er
ligeledes en lille Kunst for sig, og kan more Læseren, som en
Fægteøvelse kan fornøje Tilskueren, uden Hensyn til den Pris,
hvorom der fægtes.

Med andre Ord foruden den bredere Læseverden, der søger
Oplysning om Stoffet, gives der ogsaa et finere udviklet Publikum,
der mere tiltrækkes af Forfatterpersonligheden end af det Æmne,
den i Øjeblikket dvæler ved. Til dette henvender Supplement\-%
bindet sig.

April 1903.
\begin{flushright}G. B.\end{flushright}
\piecerule

% --- p. 4 ---
% blank
\opage{4}

% ============================================================
%  Sectional title, p. 5: `STRIDEN OM TRO OG VIDEN', one line of roman
%  capitals, centred, a little above the middle of an otherwise empty
%  page; no rule, no folio.  Indhold (PDF 9) sets it in italic:
%  `Striden om Tro og Viden 5', with the pieces under it indented.
% ============================================================
% --- p. 5 ---
\clearpage
\phantomsection
\addcontentsline{toc}{chapter}{Striden om Tro og Viden}
\begin{center}\opage{5}{\Large STRIDEN OM TRO OG VIDEN}\end{center}

% --- p. 6 ---
% blank
\opage{6}

% ============================================================
%  Samvittighed og Videnskab (printed pp. 7--10)
%  Indhold (PDF 9) reads: `Samvittighed og Videnskab 7'
%  Head as printed: `A. C. Larsen: Samvittighed og Videnskab' / `(1866)'
%    -- `A. C.' with a space and NO hyphen, confirmed at 600 dpi
%    (witness 2 reads `A.-C.').
%  Running head (pp. 8--10) reads: `Samvittighed og Videnskab'.
%  The text opens with a two-line drop capital S; set plain.
%  Dated at the end `I December 1865.', against the head's `(1866)':
%  both kept as printed.  An Efterskrift paragraph follows the date;
%  only its first word and the names in it are italic.
% ============================================================
% --- p. 7 ---
\piecehead{Samvittighed og Videnskab}{\opage{7}\emph{A. C. Larsen: Samvittighed og Videnskab}\\[2pt](1866)}
Skrifter af filosofisk Indhold egner sig i Almindelighed ikke
synderlig til Anmeldelse i et Dagblad, og vover en eller anden
begejstret Mand engang Forsøget, falder det, som Erfaring
viser, lettelig ud til et ‹videnskabeligt Digterværk› i Løjtnant v.
Buddinges Stil. Da nu lærde Bedømmelser med Rette anses
for utilstedelige, er Følgen den, at vore Blade angaaende viden\-%
skabelige Æmner vrimler af yderst korte, ubeviste og tidtnok
ubevislige Paastande, som Publikum vel har faaet en sand
Virtuositet i at nedsvælge ganske raa, men hvis Tal alligevel
intet forstandigt Menneske kan føle Lyst til at forøge. Der er
da ikke Meget, som indbyder til at anmelde et Flyveskrift, som
ikke hører Underholdningsliteraturen til; men da dette med
mere end almindelig Dygtighed behandler en Sag, som ikke blot
har Krav paa udbredt Deltagelse, men forlængst har lagt Beslag
paa en saadan, bør den Læsekres, Skriftet kan vente sig, gøres
opmærksom paa dets Fremkomst.

Det indeholder en, tildels i Meddelelsens Form, men hyp\-%
pigere forsvars- og angrebsvis fremsat teologisk Grundbetragtning,
og har som videnskabeligt Skriftemaal særlig Interesse, idet
Forfatteren besidder en filosofisk Dannelse, der bliver sjældnere
og sjældnere blandt Nutidens Teologer, og som hjælper ham
over de første Skær, paa hvilke Flertallet med ophøjet Kold\-%
sindighed sidder fast i et fælles Skibbrud.

Skønt Bogen kun paa enkelte Steder ligger i Fejde mod
Nulevende, vil den dog rimeligvis blive betragtet som et
Stridsskrift, og Hr. \emph{Larsen} vil for en Del selv bære Skylden,
eftersom han i et Forord navnlig har betonet sit angribende
% --- p. 8 ---
\opage{8}Sigte. Den Del af Bogen, der har øjeblikkelig Interesse, er
desuden den bedste, og man maa være Forfatteren taknemmelig,
fordi han har havt Mod og Vilje til at indlade sig paa en Drøf\-%
telse af de Grænsespørgsmaal angaaende det indbyrdes Forhold
mellem Tro og Viden, om hvilke der saa sjældent høres et klart
og oplysende Ord; Trangen til at faa disse Spørgsmaal besvarede
har allerede længe hos de Studerende trængt al anden Trang til
filosofisk Kundskab tilbage, og ikke én af Videnskabens indre
Opgaver ligger dem halvt saa meget paa Hjerte som denne
Grænsestrid. Vel har Adskillige slaaet sig til Ro, men nærmest
saadanne, der vilde føle sig fuldkomment tilfredsstillede ved
enhversomhelst forsonlig Løsning; det er jo menneskeligt, at
man helst vil have baade i Pose og i Sæk, og det er naturligt,
at man saa ikke tager det saa nøje med, om man faar rent
Brød i Posen eller ej. Men \emph{deres} Antal er langt større, som
efter faa Aars Forløb næsten ikke har beholdt andre Indtryk
tilbage af Professor \emph{Nielsens} aandfulde og frugtbare Lærervirk\-%
somhed end dem, der forholder sig til de nævnte Vanskeligheder,
og for hvem disse Spørgsmaal bestandig er brændende. Denne
Tilstand er hverken god eller sund, saa meget mere som hine
Indtryk aldrig er rigtig klare. I Professor Nielsens Filosofi, der
endnu bestandig udfolder sig frodigt for de Studerendes Øjne, er
nemlig dette Punkt i streng Forstand det springende Punkt, den
egenlige ‹Uro›, der gav Tankebevægelsen Fart, drev Opga\-%
verne frem og forhindrede den geniale Arbejder fra noget Øje\-%
blik at søge Hvile i stivnede Former. Men det har syntes, som
var det en Betingelse for det Heles Liv, at dette inderst æggende
Spørgsmaal bestandig blev holdt aabent, at Vunden ikke lukkede
sig, hvor Tænkeren var dybest saaret af Guden. Ikke som om
der intet Svar blev givet, men Besvarelsen er mere skikket til
at fremkalde nye Spørgsmaal end til at bringe de allerede op\-%
kastede til at forstumme.

Allerede i Professorens \emph{Propædeutik} fra 1857 indeholder de Af\-%
snit, der behandler det omtalte Forhold, en tydelig Selvmod\-%
sigelse; de lader Tro og Viden, skønt ‹absolut uensartede Prin\-%
ciper›, forsones i én Bevidsthed, eftersom Striden imellem dem
‹ikke hidrører fra det Principielle›, og denne Modsigelse er ikke
senere bleven hævet. Professor Nielsens Standpunkt er her over\-%
hovedet et saadant, at han vil kunne finde to Slags Modstandere,
for det Første dem, der forsvarer den bestaaende Teologi, for
% --- p. 9 ---
\opage{9}det Andet dem, der angriber den i dens Rod og Grund; hine
vil Prof. Nielsen tillivs som Teologiens Angriber, disse, der søger
at kuldkaste Teologien, maa stræbe at godtgøre deres Ret til at
kalde Prof. Nielsen selv for Teolog. Den Modstander, der nu
er optraadt, staar i de førstnævntes Række. Han kæmper som
Teolog for Alter og Arne, uden Nag eller Nid, men ikke uden
Blik for de Blottelser, hans overlegne Modpart frembyder.
Det lykkes ham ikke at bevise, hvad han egenlig havde sat sig
til Formaal, nemlig at Teologien endnu har Livskraft, og der
turde være gode Grunde for, at et saadant Bevis maatte glippe;
men han godtgør utvivlsomt, at Prof. Nielsen ikke fra sit Stand\-%
punkt kan kaste Vrag paa Teologien, selv om en vis bestemt
Teologi og en vis bestemt Dogmatik er sprængte.

% `Ærefrygt' as printed (witness 1 reads `Erefrygt').
Forfatteren har endnu en Fortjeneste. Den naturlige Ærefrygt
for Kierkegaards Minde har paabudt Tavshed om hans Værker;
de har faaet Lov til at tale uden at nogen summende Flok af
Ræsonnerende har forstyrret Indtrykket af den mægtige Monolog.
At dette var Ret, vil enhver Dansk forstaa; en Tysker vilde
aldrig fatte det, og Kierkegaards norske Trompeter, Hr. Dietrich\-%
son, maaske allermindst. Men ligefuldt, det var det Rette. Dog
har deraf fulgt en Ulempe, nemlig en Forstening af visse Kierke\-%
gaardske Sætninger i skintilforladelige Fortolkninger. Paa saa\-%
danne Punkter maa da en Undersøgelse af Kierkegaards Mening
sættes i Gang, og i Afhandlingen om Paradokset har Hr. Larsen
indledet en saadan.

Men naar nu disse Fortrin er fremhævede, saa er dermed
ogsaa udtalt det Bedste, der er at sige om Bogen. Forfatteren, hos
hvem man trods overvejende Selvstændighed sporer Jacobis og
Schleiermachers Indflydelse, mangler dels endnu til en vis Grad
Skole, dels i høj Grad Dialektik. Den første Mangel røber sig
f. Eks., naar han i Slutningen af sin Afhandling om Samvittig\-%
heden saa grundigt godtgør Umuligheden af at føre et sagligt
Bevis for Guds Tilværelse, at han med det samme kommer til
at godtgøre Umuligheden af overhovedet at slutte fra Virkning
til Aarsag, eller naar han i sin noget lægmandsagtige Fejde
mod Idealisterne beviser de livløse Genstandes Tilværelse af
det selvbevidste Væsens Eksistens, før dette er blevet sig be\-%
vidst. Mangelen viser sig fremfor Alt i en Sætning som denne:
‹Miraklet stadfæster netop Naturloven; sætter man Undtagelsen,
maa man ogsaa sætte Reglen›, hvori Udtrykkene \emph{Lov} og \emph{Regel},
% --- p. 10 ---
\opage{10}ganske som de to berømte Hajfiske, af Hjertenslyst sluger hin\-%
anden indbyrdes.

Af endnu større Betydning er Mangelen paa Dialektik. For\-%
fatteren taler Filosofiens Sprog, han bruger dens Betegnelser, han
tumler med dens Udtryk; men --- Røsten er ikke dens. Man mærker
i hver Linje Teologen. Forfatterens ‹Absolute› er --- undertiden for\-%
bavsende --- udialektisk opfattet, og hvorfor? Fordi dette Udtryk
i hans Mund ikke betegner Filosofiens Idé, men kun er et Navn,
et Skjul, en Maske, bag hvilken han søger at gøre Teologiens
Gudsforestilling ukendelig. Men idet nu Begrebsbetegnelserne
virker tilbage paa Tankegangen, indvikler han sig ikke sjældent
i Modsigelser; snart skal f. Eks. Mennesket ikke være væsens\-%
forskelligt fra det Absolute (det Evige, det Sande), snart skal det end
ikke kunne begribe Muligheden af det Absolutes Handlinger. Staar
nu Forklaring mod Forklaring, Filosofiens mod Teologiens, da
indtager Forfatteren følgende Stilling. Da det skorter ham paa
dialektisk Evne, bliver den filosofiske Forklaring ham uforstaae\-%
lig, Mirakelforklaringen giver sig jo end ikke ud for at kunne
forstaas, og det gælder da et Valg imellem dem. For dette
Valg gør han saa ikke Fornuften, men Hjertet ansvarligt. At
man ikke vil godkende Mirakelforklaringen, det forvirrer ham da
ikke, det forstaar han saa vel; ti det beroer paa Hjerternes
Forhærdelse. Her kaster Hr. Larsen det filosofiske Skrud til Side
og viser sig i sit teologiske Ornat. Her er da ogsaa Grænsen for
den videnskabelige Drøftelse.

% dateline, indented as a paragraph; it reads 1865 against the head's (1866).
I December 1865.

% Efterskrift: only `Efterskrift', `Dagbladet', `Rudolf Schmidt' and
% `Fædrelandet' (its genitive -s roman) are italic; the stop after
% `Efterskrift' and the comma after `Schmidt' are not clearly italic.
\emph{Efterskrift}. I en Artikel om Hr. Larsens Bog i \emph{Dagbladet}
for i Gaar erklærer Hr. \emph{Rudolf Schmidt}, at det ikke ‹for denne\-%
sinde› er hans Hensigt at opklare den ovenfor omtalte Dunkelhed
i Professor Nielsens Lære. Dette triste Budskab bringes herved
ogsaa \emph{Fædrelandet}s Læsere.
\piecerule

% ============================================================
%  Svar til Hr. Schmidt (printed pp. 11--12)
%  Indhold (PDF 9) reads: `Svar til Hr. Schmidt 11'
%  Head as printed: `Svar til Hr. Schmidt' / `(1866)'
%  Running head (p. 12) reads: `Svar til Hr. Schmidt'.
%  No drop capital here; the text opens with an ordinary indent.
%  Ends with a line of English verse in italic, set at the paragraph
%  indent, then a short rule.
% ============================================================
% --- p. 11 ---
\piecehead{Svar til Hr. Schmidt}{\opage{11}\emph{Svar til Hr. Schmidt}\\[2pt](1866)}
Foruden et mislykket Forsvar indeholdt Hr. Larsens Flyveskrift
et Angreb, der førte Krigen over i Modpartens eget Land. Ind\-%
hugget var ulige bedre end Paraden; thi en god Sag væbner
Øje og Haand. I min Lørdagsartikel har jeg stræbt at sætte
Angrebets Berettigelse i et skarpere Lys, og dette Punkt har jeg
fastholdt som Hovedpunktet; thi paa det er det, at Interessen
hviler. Kampen mod Teologien er af gammel Dato, den har
\emph{Feuerbach} og \emph{Kierkegaard} hver fra sin Side ført med en Kraft
% `er' (before `bleven'): the e is broken type, printing as an open
% hook; read `er' at 600 dpi (witness 2 reads `or').
og Følgestrenghed, som aldrig senere er bleven naaet; paa den
Kamp ligger nu ringe Vægt; men saa vist som Professor \emph{Nielsens}
Filosofi er en langt betydeligere og langt intelligentere Magt end
den teologiske Videnskab i vore Dage, saa vist har med Hensyn
til den ethvert Tvivlsmaal, der ikke er grebet ud af Luften, og
ethvert Stød, der ikke rammer Luften, en stor og væsenlig Be\-%
tydning. Men end større Opmærksomhed fortjener Sagen, naar
dens Pointe er et Midtpunkt, naar Spørgsmaalet er et Livsspørgs\-%
maal af den yderste Vigtighed, af gennemgribende Indflydelse, og
denne Sag er saaledes beskaffen. Den er det, Hr. \emph{Rudolf Schmidt}
behandler som en Ubetydelighed.

Hvis nu den Modsigelse, jeg har paapeget, virkelig var hævet
i noget af Professor Nielsens trykte Skrifter, da var jeg visselig
en større Charlatan i Filosofien, end nogen af de Flere, vort
Fædreland besidder. Men det forholder sig ikke saaledes. Hr.
Schmidt er saa god at sige mig en Artighed for filosofisk
Dannelse, der har en Henvisning til \emph{Grundideernes Logik} paa
Slæbetov. Jeg har saa lidet Brug for den sidste som for den
% --- p. 12 ---
\opage{12}første; jeg begriber end ikke, hvad Hr. Schmidt har havt i Sinde
med den. I det udkomne Afsnit af \emph{Grundideernes Logik} findes,
med Undtagelse af nogle som blotte Paastande fremsatte Sætninger
i den almindelige Indledning, Intet, aldeles Intet denne Sag vedrø\-%
rende og kunde efter Anlæggets Natur heller Intet findes. Først ved
Overgangen fra Videns Kres til Magtens kan Spørgsmaalet komme
frem, om end i virkelighedsfjernere Form, og sin egenlige Be\-%
svarelse kan det først faa ved hele Logikens Overgang til Natur\-%
filosofien.

Jeg har ikke ventet Oplysningen af Hr. Schmidt; jeg har
kun gjort Løjer med et Udtryk af ham. Jeg har en grun\-%
det Tvivl om, at en Polemik med ham vilde blive lærerig og
skal derfor ikke indlade mig paa den; ti hvad jeg ønsker er, at
lære. Men derfor vilde ogsaa en Forklaring fra Prof. Nielsens
Side være mig af den højeste Værdi, og Intet skulde glæde mig
mere, end om han vilde føle sig opfordret til at meddele en
saadan. Ti i dette Anliggende trænger vi ikke til Tavshed eller
Tilbageholdenhed, men til Ord, der gaar fra Leveren og til
Hjertet, til hensynsløs Frimodighed, til Klarhed, til Lys:

% one line of English verse, italic, at the paragraph indent; the commas
% after `Truth' and `torch' may be roman, but the line is set as one span.
\emph{Truth, like a torch, the more it's shook, it shines.}
\piecerule

% ============================================================
%  Filosofi og Teologi (printed pp. 13--16)
%  Indhold (PDF 9) reads: `Filosofi og Teologi 13'
%  Head as printed: `I Anledning af Hr. Harald Høffdings' /
%    `Stridsskrift: Filosofi og Teologi' / `(1866)'  (both lines italic)
%  Running head (pp. 14--16) reads:
%    `I Anledning af Hr. H. Høffdings Stridsskrift: Filosofi og Teologi'
%  In this piece the Google scan prints several commas as bare points
%  (tail lost); where the syntax needs a comma, the next word is lower
%  case and witness 2 (Runeberg, read from its own images) has a comma,
%  a comma is transcribed and the site is commented.
% ============================================================
% --- p. 13 ---
\piecehead{Filosofi og Teologi}{\opage{13}\emph{I Anledning af Hr. Harald Høffdings}\\
\emph{Stridsskrift: Filosofi og Teologi}\\[2pt](1866)}
Naar man gennemlæser dette Skrift for at gøre et flygtigt
Bekendtskab med Forfatterens Ansigtstræk, modtager man et
næsten uhyggeligt Indtryk; ti med enkelte og dobbelte Gaase\-%
% `Gaase-/øjne': the hyphen falls at the line end; read as a break, not
% a compound hyphen (no other occurrence to decide by).
øjne i hundredevis --- undertiden med et Antal af tredive paa en
eneste Side --- stirrer Forfatteren lig en Argus af Gaaseslægten
En tungsindig i Møde fra Bogens Papir. Hos ham spiller nem\-%
lig Citaterne den samme Rolle, som den selvtilvirkede Tekst hos
andre Forfattere.

Idet Hr. \emph{Høffding} gaar ud fra, at den filosofisk-teologiske
Strid i Virkeligheden allerede er afgjort, gælder det for ham
kun om at gøre det klart, at Prof. \emph{Nielsen} er den, som har sej\-%
ret, og i denne Hensigt er det, at han giver en tung og pakket,
efter Læserens forkjellige Standpunkt enten utilstrækkelig eller
unødvendig, Genfremstilling af Kampens Historie. Mod Hr. \emph{Lar\-%
sens} Flyveskrift:
% PRINTER'S ERROR: the opening quotation mark is printed turned, as a
% closing mark (›...›); compared at 600 dpi against ‹skønt› on p. 14.
% Transcribed as printed; the ‹/› count of this page is off by 2.
›Samvittighed og Videnskab› fremfører han dernæst
en Del Bevisgrunde, mest saadanne, man allerede én Gang har
set paa Tryk, og retter til Slutning mod Undertegnedes fore\-%
gaaende Artikler i \emph{Fædrelandet} nogle Indsigelser, der udgør Bogens
svageste, men maaske ogsaa nyeste Parti. For disse Indsigelser
er jeg Forfatteren meget taknemmelig; thi de godtgør tilfulde,
hvad jeg tilforn har udtalt, nemlig hvor stærkt der trænges til
en Redegørelse fra Prof. Nielsens Side.

Jeg har erklæret det for en Modsigelse, at Prof. Nielsen
paa én Gang kalder Tro og Viden absolut uensartede Principer
% --- p. 14 ---
\opage{14}og samtidig taler om deres Forsoning i den menneskelige Be\-%
vidsthed. Hr. Høffding lærer mig, at efter Prof. Nielsens Anskuelse
kan denne Forsoning ské, ikke \emph{skønt}, men \emph{fordi} disse Principer
er absolut uensartede. Dette er temmelig langt fra at være no\-%
gen Hemmelighed; men Modsigelsen er ligefuldt betegnet ved
mit ‹skønt›. Er det tænkeligt, at to absolut uensartede Prin\-%
ciper kan forliges i samme Bevidsthed? Lader det sig godt\-%
gøre, at Tro og Viden er slige? Jeg besvarer disse Spørgsmaal
benegtende. Det turde dog være Tvivl underkastet, om der i
den menneskelige Aand kunde raade to absolut uensartede Prin\-%
ciper, uden at der indtraadte en Spaltning, hvorved hele Tilvæ\-%
relsen revnede, og at Prof. Nielsen selv ikke har taget det saa
strengt med denne absolute Uensartethed, beviser den Omstæn\-%
dighed jo bedst, at han taler om Strid og Forsoning mellem
hine Magter; ti absolut uensartede Magter kan hverken føre
Krig eller slutte Fred med hinanden. Men lad os engang
sætte, at vi her har to absolut uensartede Principer for os, ser
Hr. Høffding saa ikke, hvor besynderlig al denne Tale bliver
om, at man skal sætte Troen øverst og Videnskaben nederst,
men ikke omvendt; ser han ikke, at ethvert Spørgsmaal om,
hvad der skal sættes højest og lavest, i saa Fald vil blive endnu
mere komisk, end \emph{Poul Møllers} gamle: Hvad er højest, en Pop\-%
pelpil eller et Tordenskrald? Og dog er maaske mangen stakkels
Studiosus kommen galt afsted, naar han paa Eksaminatoriet sva\-%
rede forkert og i sin Uskyldighed vendte op og ned paa Tro og
Viden.

For at lette Forstaaelsen giver Hr. Høffding imidlertid
endnu et Vink. Etiken, den filosofiske Etik, ‹danner en Over\-%
gang mellem Filosofi og Religion›. I og for sig selv betragtet
kan denne Sætning være meget fortræffelig; men i Hr. Høffdings
Mund tager den sig rigtignok uforlignelig ud. Ti vistnok har
Moralen altid været anset for en beundringsværdig Magt, men
at den kan danne Overgangen mellem to absolut uensartede
Krese, det har neppe Nogen tiltroet den. Mere end til at
undersøge, hvorledes dette skulde gaa til, fristes man til at
spørge, hvad der har bevæget Hr. Høffding til at sende denne
Snedighed ud i Verden. Svaret er ikke vanskeligt.

For Enhver, der med et opmærksomt Øje har fulgt Gangen
i Professor \emph{Nielsens} rige Forfattervirksomhed i det sidste Tiaar,
er dét efterhaanden bleven klart, at han, uden at tage noget ud%
% --- p. 15 ---
\opage{15}talt Ord tilbage, paa det omstridte Punkt har, om ikke skiftet,
saa dog ændret og lempet sin tidligere Anskuelse. Jeg indbilder
mig langt fra hermed at udtale noget Nedsættende. En Ved\-%
kendelse af, at man har forandret sin Overbevisning fra Grun\-%
den af, det kan være et ædelt, ja beundringsværdigt Træk hos
en fremragende Personlighed, og som et saadant viste denne
Handling sig, da i sin Tid Professor Nielsen sluttede sig til
Kierkegaard. Tillempende at udvikle den samme Grundtanke,
% `Noget,': the Google scan shows a bare point (600 dpi); both
% witnesses read a comma and `ikke' follows in lower case.
det er, om ellers Noget, ikke blot en ærlig, men en hæderlig
Sag. Det Sidste har efter min Opfattelse Professor Nielsen gjort.
1857 erklærede han, at mellem det Evangeliske og det Mytiske
var der i Videnskabens Øjne ingen Forskel; men 1864 lempede
han dette Udtryk saaledes, at den religiøse Etiks Overlegenhed
over den rationelle forhindrede Videnskaben fra at erklære Evan\-%
geliet for Myte. Det Sidste er vel en haard Tale; det lader
sig vanskeligt forstaa, at der kan gives to Slags Moral, og at
nogen anden Moral kan være højere, end den rationelle, ɔ: den
paa fornuftig Erkendelse af Menneskets Væsen begrundede Mo\-%
ral. Men lige meget, det er tydeligt, at her er sket en Lem\-%
pelse; ti hermed er i Grunden den absolute Uensartethed af Tro
og Viden opgivet, skønt ikke tilbagekaldt.

For Sligt har nu Hr. Høffding intet Blik; en Udvikling, der
trods enkelte Ytringers tilsyneladende Protest gaar i den Ret\-%
ning at løfte Videnskaben fra det underordnede Forhold, hvori
den tidligere var sat, kan han ikke se; men frejdig slaar han
løsrevne Sætninger fra de forskelligste Bøger i Bunke og stabler
den ene Modsigelse ovenpaa den anden. Han fastslaar, at hver\-%
ken det Rette, Sande eller Gode vilde kunne holde sig i det virke\-%
lige Liv, hvis det ikke støttedes af ‹Religionen›; men ligefuldt
er han overbevist om, at der gives en filosofisk, altsaa rationel,
Moral, der --- under den nysnævnte Forudsætning --- i det mindste
vil have den, mindre fornuftige, Egenskab, at der ikke kan leves
efter den. Han hævder med yderste Iver, at Samvittighedsfor\-%
dringer ikke maa gøres gældende i Videnskaben, men ligefuldt
bebrejder han med største Ugenerthed Fritænkernes Videnskab,
at den er en ‹vantro› Videnskab. Hr. Høffding er ingen Ven af
Følgerigtighed.

Paa Skriftets sidste Side retter Forfatteren den Indvending
imod mig, at, hvis jeg ikke har fundet Bidrag til Svar paa mit
Spørgsmaal i det hidtil Udkomne af \emph{Grundideernes Logik}, saa burde
% --- p. 16 ---
\opage{16}jeg ikke have indskrænket mig til at ønske Oplysning af Pro\-%
fessor Nielsen, men maatte have fundet Fejl i det Givne. Denne
Indsigelse var let at forudse og er let at gendrive. Vel er det
mig ufatteligt, hvorledes Professor Nielsen kan lade Miraklets
% `strande': the t is broken type (prints like `l'); witness 1 reads `sirande'.
Mulighed erholde Plads i sit System uden at strande i sluttelig
Splittethed; men i den udkomne Part af \emph{Logiken}, er der jo heller
ikke gjort nogetsomhelst Forsøg derpaa. Jeg indser heller ikke,
% PRINTER'S ERROR: `bebuder.' -- a point where the syntax wants a comma (600 dpi: round point, no tail), as witness 1
% reads; witness 2 reads a comma and `uden' follows in lower case.  Transcribed as printed (verification 2026-09-22).
hvorledes Udviklingen skal kunne give, hvad Fortalen bebuder.
uden at den højere Enhed af Viden og Magt, af Idealitet og
Realitet, som er Værkets Grundtanke, bliver sprængt, nærmest
derved, at Magtens Idé forvandler sig til Forestillingen om Al\-%
magt i teologisk Forstand. Dertil har jeg hentydet i en tid\-%
ligere Artikel.

Der gives her til Lands, som ogsaa andet Steds, f. Eks. i
Frankrig, visse Forfattere, der bedst kan lignes ved de bekendte
store Blærer, der i sig selv er flade, slappe og indholdsløse
Ting, men som ved Hjælp af Inspirationen udenfra svulmer saa\-%
ledes op, at de indtager en ganske anselig Plads. Det bør siges
% PRINTER'S ERROR: `Heffdings' -- the second letter is a wrong-fount
% (italic-looking) e, not ø (600 dpi); both witnesses read `Heffdings'.
til Hr. Heffdings Ære, at til denne Klasse hører ikke han. Han
har øjensynlig en redelig Bestræbelse og et alvorligt Sind. Men
disse Fortrin har ikke forhindret, at han i sit Forhold til Pro\-%
fessor Nielsen ved denne Lejlighed har vist sig som et altfor
aabenmundet Barn.
% The closing rule is printed faint and broken in the scan.
\piecerule

% ============================================================
%  En sidste Replik (printed pp. 17--21)
%  Indhold (PDF 9) reads: `En sidste Replik 17'
%  Head as printed: `En sidste Replik' / `(1866)'
%  Running head (pp. 18--21) reads: `En sidste Replik'.
%  p. 17 foot carries the signature `G. Brandes: Samlede Skrifter. XIII.'
%  and `2', p. 19 the signature `2*' -- not transcribed.
%  p. 21 is a short page (six lines) with running head and folio, ending
%  with a short rule.
% ============================================================
% --- p. 17 ---
\piecehead{En sidste Replik}{\opage{17}\emph{En sidste Replik}\\[2pt](1866)}
Enhver Teatergænger vil have gjort den Erfaring, at selv den
daarlige Skuespiller kan høste Klap og Bifald, naar han som
Uskyldighedens eller Troens Ridder udtaler en ædel Grundsætning,
og at Publikum stundom næsten med Mishag hører Skuespilleren
forfægte en Sag, som det anser for slet. Hvis min ærede Mod\-%
stander, Hr. \emph{Høffding}, kunde faa Rollerne saaledes fordelte imel\-%
% PRINTER'S ERROR: `at kan kom' for `at han kom'; both witnesses read
% `kan'.  Transcribed as printed.
lem os, at kan kom til at spille en Lysets Engel og jeg en Mør\-%
kets Dæmon, saa havde hans Sag unegtelig sejret, og den ud\-%
vortes Triumf var ham vis. En saadan Vending søger han at
give Forholdet, idet han beskylder mig for ikke at ville være
Part, men Dommer, for ikke aabent og ærligt, men kun ad Om\-%
veje at vedkende mig en fritænkersk Opfattelse, da dog ‹Ingen
tør deltage i Striden uden at melde, hvad han fører i sit Skjold›.

Men denne Vending skal ikke lykkes ham; det er ikke Vantro
mod Tro, snarere er det Klarhed mod Taage, der her staar lige
overfor hinanden. Ti hvad er Undersøgelsens Genstand, hvorom
spørges der ene og alene i denne Sag? Ganske vist ikke om \emph{min}
Filosofi, saa lidt som f. Eks. om Hr. Høffdings. Her er kun et
eneste afgørende Punkt, som det gælder om at holde fast, og
det er dette: Er Professor \emph{Nielsen} i en Modsigelse, naar han
kalder Tro og Viden absolut uensartede Principer, og dog lader
dem forsones i den menneskelige Bevidsthed, saaledes at det
ene sættes øverst, det andet nederst, eller er han det ikke? Han
er det; ti enten er Tro og Viden ikke absolut uensartede, eller
de spalter Bevidstheden helt; enten er de ikke absolut uens\-%
artede, eller der kan ikke tales om et Højest og Lavest. Det
% --- p. 18 ---
\opage{18}nytter ikke som Hr. Høffding at indvende, at det jo er det
samme Menneske, der har et Forhold til begge de to absolut
uensartede Principer; ti Spørgsmaalet er netop det, om ikke
Menneskets Bevidsthed maa splittes ulægeligt derved. Det hjæl\-%
per ikke heller, som Hr. Høffding, at sige, at derfor rettes til
Enhver personlig det Spørgsmaal, hvorledes han her vil træffe
sit Valg. Ti hvorpaa gaar Valget ud? Hvis der vælges mellem
Tro og Viden paa den ene Side og Viden paa den anden, da
vælger den, som griber det Første, jo Spaltningen; vælges der
derimod mellem Tro paa den ene Side og Viden paa den an\-%
den, da er jo Forsoningen i samme Bevidsthed opgivet og kun
Dualismen, kun Adsplittelsen bliver, om den end forlægges fra
det enkelte Menneske til Menneskeheden. Det nytter overhove\-%
det ikke som Per Degn at besvare simple og tydelige Spørgs\-%
maal med højst gaadefuld og uforstaaelig Tale.\fnmark{*)}{For at ikke Hr. Høffding paany skal sende mig mine egne Vendin\-%
ger tilbage under de Tegn, jeg ikke vil nævne, og tilmed her se bevist, at
jeg behandler Professor Nielsen som ‹Delinkvent›, tilføjer jeg, at af denne
Lignelse lige saa lidet følger, at jeg i det Hele er nogen Erasmus Monta\-%
% PRINTER'S DEFECT: the quotation opened at `‹at slaa' is never closed;
% only the inner ‹Probe ...!› closes.  Transcribed as printed (‹ +1).
nus, som at jeg har bedrevet den Udaad som hin ‹at slaa Rasmus Niel\-%
sen 3 til 4 Gange bag i Nakken med knyttede Næver og derhos altid raabe:
‹\emph{Probe Majoren, probe Majoren!}› (1. Akt 6. Scene)}

Jeg anser mig ikke, som Hr. Høffding giver sine Læ\-%
sere at forstaa, for ufejlbar; jeg anser ikke engang Profes\-%
sor Nielsen derfor, hvad Hr. Høffding uden Tvivl meget let\-%
tere havde tilgivet mig; men er jeg end ikke saa ubeskeden at
antage det for umuligt, at jeg kan fejle, saa er jeg dog endnu
langt mindre saa taabelig, at jeg uden en klippefast Tillid til
Sandheden af min Sag skulde have sat Pen til Papiret mod en
saa højtbegavet og mig saa overlegen Mand som Professor Niel\-%
sen. Hvad jeg fører i mit Skjold, er altsaa dette: en Overbevis\-%
ning, fuld og fast, om at enhver Dualisme (Tvehed) er umulig,
og at som Følge deraf Professor Nielsens Lære paa dette
Punkt er uholdbar. Hvad Filosofi jeg forøvrigt har, eller om
jeg overhovedet har nogen, det er dette Spørgsmaal gan\-%
ske lige saa uvedkommende, som hvad jeg mener om Tyrker\-%
nes Affærer.

Jeg har styrket min Bevisførelse, idet jeg godtgør, at Profes\-%
sor Nielsen selv ikke har kunnet fastholde den absolute Uens\-%
artethed. Hvis denne virkelig fandt Sted, da maatte enhversom%
% --- p. 19 ---
\opage{19}helst Form af religiøs Tro kunne forenes med Professor Niel\-%
sens Filosofi. Men spørger man f. Eks.: kan man paa en Gang
være Katolik og Tilhænger af Professor Nielsens Lære? saa er
Svaret kun i saa Fald: Ja, ifald Tro og Viden er hinanden al\-%
deles uvedkommende, men: Nej, ifald der skal tages Hensyn til
Katolicismens Etik; ti det er dog vel utvivlsomt, at denne ikke
er den rationelle overlegen. Idet Professor Nielsen lader Over\-%
legenheden af Religionens Etik være det Afgørende, er hans
første Standpunkt opgivet. Da Hr. Høffding imidlertid tror at
kunne forsvare Sætningen om den dobbelte Etik som Overgang
mellem Filosofi og Religion, skal jeg gerne følge hans Tankers
underlige Gang.

Der er, siger han, en Modsætning mellem Videnskaben og
Livet. Idet det Godes Idé i Gerning skal virkeliggøres, træder
denne Modsætning saaledes frem, at den rationelle Etik viser
sin Utilstrækkelighed, og Nødvendigheden ses --- ikke af den
religiøse Etik, som man foreløbig maatte vente, men --- af ‹Re\-%
ligionen›. Nu spørger jeg: hvad forstaas da ved den religiøse
Etik? Er den en Videnskab, eller er den blot et andet Navn
for det Positivt-Religiøse? Er den det Sidste, hvorfor da den
forstyrrende Benævnelse? Men er den det Første, saa er Sagen
simpel og Forvirringen klar. Ti da maa jo i Kraft af den selv\-%
samme Modsætning mellem Videnskab og Liv ogsaa denne Etik
blive utilstrækkelig, saa snart som den skal bruges i Livet, og
der maa atter tyes til noget Højere. Men, ikke at tale om, at
den rationelle Etiks Betydning, som dog skulde hævdes, ved
denne Tankegang indskrænkes til mindre end Intet, saa faar vi
her en Række af Videnskabeligt, Overvidenskabeligt, Overover\-%
videnskabeligt osv., der ikke taler til Fordel for dette Forsvars
Logik. Hvis det at være et sædeligt Menneske indeholdt saa\-%
danne logiske Forviklinger, saa ser jeg vel, at man med Hr.
Høffding maatte spørge: Kan man paa en Gang være Etiker og
Videnskabsmand? Men ellers er dette Spørgsmaal ikke let at
forstaa. --- Nej, det staar fast: gør Etiken Overgangen fra Filo\-%
sofi til Religion, saa er disses absolute Uensartethed opgivet, og
det bedste Bevis for denne Læres Uholdbarhed er det, at
Professor Nielsen ikke blot har opgivet den her, men, som
jeg har vist det, ogsaa paa andre Punkter. Saaledes,
naar der gives Miraklets Mulighed Plads i \emph{Grundideernes}
% --- p. 20 ---
\opage{20}\emph{Logik}.\fnmark{*)}{Det blev mig først betydet, at Problemets
% `Løsning': the slash of ø has not printed in the Google scan (the
% footnote type is faint here); both witnesses read `Løsning'.
Løsning var givet i
den udkomne Del af \emph{Logiken}. Jeg svarede, at den ikke fandtes der og
efter Værkets Anlæg ikke kunde findes. De Samme, der tænker sig
Tro og Viden afspærrede i to forskellige Rum i Bevidstheden, beskyldte
mig da for at hylde en Rubrikfilosofi. Jeg svarede, at jeg forstod Vær\-%
kets Grundtanke saaledes, at Miraklets Mulighed ikke lod sig forberede
eller udlede, men kun kunde fremkomme ved et Spring. Nu negter Hr.
Høffding sig ikke den Bebrejdelse mod mig, at jeg ikke har paavist den
Fejl i det allerede Udkomne, som jeg selv erklærer der ikke findes.}
Ti Forskellen mellem Filosofiens Absolute og Teolo\-%
% `den,': the Google scan shows a bare point; both witnesses read a comma.
giens almægtige Gud er den, at denne ikke som det Absolute
er betinget af sine egne Betingelser; men denne Forskel hæ\-%
ver Logiken, idet den sætter Miraklets Mulighed. I Logiken
fremkommer da den samme Fordobling som i Etiken. Lige\-%
som det ikke lader sig forstaa, hvorledes den religiøse Etik
kan undlade at faa tilbagegaaende og tilintetgørende Virkning
paa den rationelle, saaledes lader det sig ikke tænke, at den
ubetingede Personlighed i Kraft af sin Selvbestemmelse kan løfte
sig op over sine egne Selvbestemmelser, uden at disses Be\-%
tydning tilintetgøres, det vil sige: uden at Naturens og Aan\-%
dens Love, som vi kender dem, bliver til Intet, til Skin. Jeg
forstaar, at man med Kierkegaard kan støtte sin paradokse
Tro ved at forkaste Videnskaben; men jeg fatter ikke, at man,
som Professor Nielsen, kan opgive Tvivlen om Videnskab og
dog beholde hin Tro, i det mindste ikke uden at man nødes
til at gøre de Samvittighedskrav gældende mod Fritænkeren,
som man ikke vil taale, at Teologen retter til En selv. Denne
Udtalelse kan maaske forklare Hr. Høffding, hvad han saa
meget interesserer sig for at vide, nemlig hvad det er, jeg
% PRINTER'S DEFECT: `Menne-/sker' is broken at the line end with no
% hyphen printed (both witnesses read `Menne sker').  Joined here.
anerkender hos Kierkegaard. Kun meget uselvstændige Menne%
sker forstaar dog ved Anerkendelse ikke Andet end en slet og
ret Vedhængen gennem tykt og tyndt. Hvad jeg beundrer hos
Kierkegaard er blandt meget Andet hans Overensstemmelse med
sig selv.

Dette er mit sidste Ord til Hr. Høffding. Nu er vel Intet
mig fjernere end af Professor Nielsens Tavshed at slutte, at
han vil godkende, hvad der af mine Modstandere er gjort gæl\-%
dende imod mig, eller at disses Indlæg virkelig skulde inde\-%
holde den rette Fortolkning af hans Anskuelse; men jeg tror,
at det havde været til Gavn for Sagen, om han havde tilkende%
% --- p. 21 ---
\opage{21}givet dette. Et haandgribeligt Resultat kunde en Tankestrid som
denne ikke bringe; men det er mit Haab, at har mine Udtalel\-%
ser ikke givet andet Udbytte, saa har de dog maaske gjort en
og anden i Stilhed Stræbende, der har Mod til at dømme, Vilje
til at dømme ærligt, Forstand til at dømme rigtigt, opmærksom
paa, hvor Vanskeligheden ligger.
\piecerule

% ============================================================
%  Om Selvansvarlighed (printed pp. 22--25)
%  Indhold (PDF 9) reads: `Om Selvansvarlighed 22'
%  Head as printed: `Om Selvansvarlighed' / `(1866)'
%  Running head (pp. 23--25) reads: `Om Selvansvarlighed'.
%  No drop capital; the text opens with an ordinary indent.
%  Ends with a small centred couplet from Peder Paars (roman, set
%  smaller), then a short rule.
% ============================================================
% --- p. 22 ---
\piecehead{Om Selvansvarlighed}{\opage{22}\emph{Om Selvansvarlighed}\\[2pt](1866)}
Forf. til ‹Julen hos Elise› har paa Gads Forlag udgivet et
Skrift ‹\emph{Om Selvansvarlighed}›, hvilket som Eksempel kan tjene
% `Enten---Eller': the dash is printed without spaces, inside the italic.
til Oplysning af følgende Sætning i ‹\emph{Enten---Eller}›: ‹Der
gives en Ræsonnementspassiar, der i sin Uendelighed staar i
samme Forhold til Resultatet som de uoverskuelige ægyptiske
Kongerækker til det historiske Udbytte›. Ti disse Ord angiver
det rette Synspunkt for en Bog, der som denne er uden Fynd
i sit Udtryk og uden Gang i sine Tanker, benløs i sin Bygning,
leddeløs i sin Sammenføjning og blodløs i sin Udførelse.

Forfatteren har stillet sig det priselige Formaal at være almen\-%
fattelig, men han har ikke indset, at Fattelighed opnaas meget bedre
ved en Tankens Tugt, der udelukker taaget Bredde, uvedkom\-%
mende Indskud, Gentagelser og Trivialitet, end ved nogen af
disse lidet glimrende Laster selv; han har fremdeles fulgt den
priselige Plan at lade alle tidligere Arbejder om dette Æmne
uomtalte; men det synes næsten, som om han tillige selv enten
ikke har læst eller ikke forstaaet dem; thi ellers er det ube\-%
gribeligt, at han end ikke har været i Stand til med Bestemthed
at \emph{stille} sit Problem, selv om han hverken formaaede at løse
det eller at hugge det over. Hans Skrift hører da saa lidet til
dem, der klarer den store Mængde en Videnskabs Resultater,
som til dem, der fører Literaturen videre; det burde ikke være
udgivet som selvstændig Bog; det kunde i forkortet Form have
hævdet sig en Plads i et af vore to teologiske Tidsskrifter.

Bogens Grundstemning er en dyb Fornemmelse af vore
Evners Ufuldkommenhed til at løse Tilværelsens Gaader. For%
% --- p. 23 ---
\opage{23}fatteren begynder altid med teologisk at hemmelighedsfylde Forhol\-%
% A raised speck stands in the space between `Om' and `Samvittig-'
% (spacing material); not transcribed.
dene for saa bagefter at erklære dem for uforklarlige. Om \emph{Samvittig\-%
heden} siges der f. Eks., at skønt den maa staa i en vedvarende
% PRINTER'S ERROR: the quotation opened at `aldeles' is closed with a
% turned mark, printed `vedligeholdes‹;' -- transcribed as printed
% (‹/› balance +1 from here).
Forbindelse med det altstyrende Princip, er det ‹\emph{aldeles ufor\-%
klarligt}, hvorledes Forbindelsen finder Sted og vedligeholdes‹; om
\emph{Livet} hedder det, at Livskraften (!) giver sig kun tilkende, naar den
er forbunden (!) med en legemlig Tilværelse, ‹men hvorledes Livs\-%
kraften forbinder sig med det Legemlige, saaledes at dette kan
producere sig i et nyt Individ, er det \emph{Uudgrundelige}›. Somme\-%
tider forudsætter han Grader af Forstaaelighed indenfor dette
Uudgrundelige, som naar han i følgende Sætning skifter Ret
imellem Teisme og Panteisme: ‹Er et levende, fornuftigt, per\-%
sonligt Væsens første Oprindelse og hele Virken som almægtig
Gud \emph{ufattelig}, da er en upersonlig Grundkrafts Oprindelse og
Produktionskraft det ikke mindre; men Verdenslivet, dets Øjemed
og Maal, bliver \emph{dog mere forstaaeligt}, naar det betragtes som
udgaaet fra en personlig Gud›. Sommetider opgiver han der\-%
imod aldeles Ævred som i følgende modløse, naive, men i For\-%
fatterens Mund naturlige Erklæring: ‹vi tilstaar, at vi aldrig
kan vinde nogen Klarhed i vor Tanke om de Ting, der ikke
lader sig fremstille for vore Sanser›.

For at vinde Fodfæste i dette Mørke gaar Forfatteren nu
snart ud fra teologiske Forudsætninger som den om en op\-%
rindelig Sædelighedstilstand, der senere er gaaet tabt, eller
som den om Arvesynden, snart appellerer han som første
% The colon after `Følelse' is clearly italic (600 dpi).
Forudsætning til vor \emph{Følelse:} den positive Religion bekræftes
af ‹vor egen indre Følelse› eller af ‹den renere Menneske\-%
naturs uvilkaarlige Følelser›, og en modsat Anskuelse, som
rigtignok ‹vel kan have Noget for sig i, hvad vi paa sanselig
Maade erfare om den Udviklingsgang, Verden har gennemgaaet›
[altsaa i det, hvorover vi efter Forfatterens Mening ene kan
komme til Klarhed], gendrives derved, at ‹der ligger Intet i den,
som kan tiltale den menneskelige Følelse›. Naar først denne
% `vé' with acute, as printed.
Rettesnor for Sandhed bliver antagen, vé da Læren om Infu\-%
sionsdyr og Trikiner; thi sandelig disse Dyr har Intet, der
tiltaler den menneskelige Følelse.

Paa det nu angivne Grundlag bestemmes Forholdet mellem
Tro og Viden saaledes: ‹Den ene er vunden ad praktisk Vej,
gennem Erfaring, ved Betragtning af det Synlige, som kan iagt\-%
tages af Alle; derfor kan Forestillingen derom ogsaa lade sig
% --- p. 24 ---
\opage{24}konferere med og lade sig berigtige efter den Forestilling, Andre
har dannet sig derom, og naa en saadan Sikkerhed, at vi kan
betegne den som \emph{Viden}. Forestillingen om det højere (!) Aande\-%
lige bliver derimod af mere (!) subjektiv Natur, fordi den ikke
har noget Synligt at knytte sig til, som er iagttageligt for Alle
% The stop after `Tro' is not clearly italic; kept roman.
paa samme Maade; vi betegner den som \emph{Tro}.› Efter denne smukke
og nøjagtige Begrebsbestemmelse bliver blandt Andet Logik og
Generalbas Tro; men med en ung, moderne Dames hele
elskværdige Ugenerthed er Forfatteren vis paa at blive forstaaet,
naar han ophænger Troen ved det rummelige Komparativ,
man kalder ‹det Højere›.

Forfatteren er lige stor som Stilist og som Filosof. I første
Egenskab ser vi ham, naar han lærer, at ‹mange Lidenskaber og
Affekter kun kan fremkomme og bestaa ved \emph{legemlig Under\-%
støttelse}, som Hidsighed og Hang til Udsvævelser› (Krop skal
der til), eller naar han taler om, hvorledes det menneskelige
Legeme ‹udvider sig fra en højst ubetydelig Begyndelse stadig i
større og større Omfang›. Som Tænker hører vi ham i følgende
kostelige Udledning af Ærgerrigheden: ‹Det skønne Legeme
vækker Beundring, men kun det menneskelige Legeme kan frem\-%
kalde Ære; thi det er alene Aanden, der, som Besidder af dette
Legeme, bliver Genstand derfor. Det er fornemmelig det aande\-%
lige Udtryk, der nærmest giver sig tilkende i Ansigtstrækkene,
som kan fremkalde Ære. Men Udtrykket er usikkert og kan
ikke hos Alle gennemtrænge Legemligheden, derfor [ergo] maa
Mennesket søge at tilvinde sig Ære ved sine Handlinger›. Hol\-%
bergs Henrik har aldrig paa vittigere Maade ført Bevis.

I Bogens syvende Afsnit: ‹Forholdet mellem den menneske\-%
lige Frihed og Guds Styrelse› har Forfatteren sammentrængt sin
% PRINTER'S DEFECT: no full stop after `Forvirring' (600 dpi: a wide
% space and nothing in it); transcribed as printed.
Forsknings Resultater og sin hele Forvirring Dette Afsnit ret\-%
færdiggør Guddommen. Ræsonnementet er følgende: Den Styrelse,
% `naaer' as printed.
som naaer sit Maal under Bevarelse af Friheden for de Personer,
hvilke den benytter som Midler, maa betragtes som den fuldkom\-%
neste. Følgelig strider Menneskets Frihed ikke mod Styrelsens Fuld\-%
kommenhed. Da Udviklingen i Kraft af Frihed nu imidlertid
ikke i alle Enkeltheder kan være overensstemmende med den
guddommelige Vilje, opstaar det Spørgsmaal, om det ikke var
Styrelsen værdigst at ophæve Frihedens slette og tidt for den
Uskyldige ulykkebringende Virkninger ved en overordenlig Ind\-%
griben, ved Miraklet. Forfatteren benegter dette. Hvorfor?
% --- p. 25 ---
\opage{25}Hans Slutningsmaade er værd at høre: ‹Hvis det Skud, der
rettedes mod den Uskyldiges Bryst, prellede af fra dette uden
at skade ham, da vilde Følgen kunne staa som et Vidnesbyrd
om hans Uskyldighed, men den vilde ikke kaste noget Lys over
den udøvede Handling; denne vilde kunne anstilles som en
Prøve paa dens Natur, imod hvem den rettedes; den vilde ikke
virke til en bedre Erkendelse af det Gode eller Slette i den
Hensigt, som ledede til Handling.› Hvorledes en saadan Er\-%
kendelse af Hensigtens moralske Beskaffenhed lader sig udlede
af Udfaldet, overlader vi Forfatteren at føre Beviset for, og vi
retter kun for at være sikre paa at forstaa ham det Spørgsmaal
til ham: Gives der altsaa intet Mirakel, intet Overnaturligt?

Man maa jo dog forbavses ved, at en saa velopdragen teologisk
Filosof som Forfatteren synes at ville svare med en ligefrem
Benegtelse. Men man berolige sig, det gør han ikke heller,
han danner sig i Filosofien efter de bedste Mønstre. Han
mener med Salomon Goldkalb ‹vieles for und bag, und ja und
nej tillige›: ‹Naar vi ikke ser nogen Grund til at antage saa\-%
danne Brud paa den naturlige Orden og Lov som nødvendige
for Guds Verdensplans Fuldendelse, saa udelukker dette dog
ikke en guddommelig Virken, der for Mennesket maa staa som
overnaturlig. En saadan maa der være; thi uden denne vilde
den i Skabningen nedlagte Kraft ikke kunne vedligeholdes (!).
Guds Virken som Verdens Opholder er for saa vidt en over\-%
naturlig Virken, som den er \emph{ufattelig} for de menneskelige Evner,
men den bliver ligefuldt naturlig, naar den ytrer sig under
Fastholdelsen af bestemte Love gennem en \emph{højere} Natur›. Lur
Forfatteren! Der fik han paa én Gang Brug for sit \emph{Ufattelige} og
sit \emph{Højere}, ja fik begge slaaede sammen til en højere Ufatte\-%
lighed.

% `teologisk-filosofisk': a compound hyphen that falls at the line end;
% kept.
Som et blandt mange Vidnesbyrd om, i hvilket teologisk-%
filosofisk Uføre vi her i Danmark for Øjeblikket sidder fast, er
denne Bog ikke uden al Interesse. Men iøvrigt véd vi intet
bedre Raad at henvende til dens Forfatter end det, som inde\-%
holdes i dette gamle Dobbeltvers af \emph{Peder Paars}:

% The Holberg couplet: roman, set smaller, centred.
\begin{center}\small
I Fingren stikke maa i Tiden, lugte til,\\
hvad Tid I lever i, om mig I lyde vil.
\end{center}
\piecerule

% ============================================================
%  Nationalitetsideens Gyldighed (printed pp. 26--27)
%  Indhold (PDF 9) reads: `Nationalitetsideens Gyldighed 26'
%  Head as printed: `Claudius Wilckens: Nationalitetsideens Gyldighed'
%    (one line) / `(1867)'
%  Running head (p. 27) reads: `Nationalitetsideens Gyldighed'.
%  No drop capital; the text opens with an ordinary indent.
% ============================================================
% --- p. 26 ---
\piecehead{Nationalitetsideens Gyldighed}{\opage{26}\emph{Claudius Wilckens: Nationalitetsideens Gyldighed}\\[2pt](1867)}
Imod en Afhandling af Provst Kofoed-Hansen: \emph{Nationalitet
og Kristendom} fremhæver Forfatteren, at Spørgsmaalet om
Nationalitetsideens Berettigelse ligeoverfor Kristendommen kun
er en enkelt Side af Spørgsmaalet om den menneskelige Kulturs
Berettigelse overhovedet, saa at man ikke kan negte hin Idé
Betydning og Gyldighed uden at bryde Staven over den hele
Kultur. For saa vidt turde Hr. Wilckens have fuldkommen Ret.

Men naar han nu gaar over til at fremsætte sin egen Opfattelse
af Forholdet mellem Kristendom og Humanitet, da er hans
Synsmaade endnu mere uklar og modsigende end hans Mod\-%
standers. Fra Filosofiens Standpunkt forkynder han --- som
det synes uden ret at kende Betydningen af det Ord, han bruger
--- en ubetinget \emph{Dualisme} af guddommelig Historie og menneskelig
Historie, af religiøs Udvikling og menneskelig Udvikling, af Forsyn
% Broken type: the second `s' of `Verdensstyrelse' is damaged (reads
% `Verdensstyrel.se'); transcribed as the word.
og Verdensstyrelse. Men Dualismen (Grundspaltningen) er den viden\-%
skabelige Tænknings Blusel, som Tænkeren selv dog skulde være
den Sidste til at stille til Skue. Skønt Forfatteren selv erklærer,
at det evige Forsyn \emph{satte} Verdensstyrelsen til at raade over
Folkenes Skæbne, fremhæver han med spærrede Typer, at Ver\-%
denshistorien, ledet af Verdensstyrelsen, \emph{Intet, slet Intet} har at
gøre med Forsynet, og dette \emph{Intet, slet Intet} at gøre med hin.

Dette er allerede at gaa meget vidt; men endnu højere spænder
% `Kofoed-Hansen': compound hyphen at the line end; kept.
Forfatteren dog Buen, naar han, bebrejdende Provst Kofoed-%
Hansen, at han mangler endog enhver Anelse om den nævnte
mærkværdige Dualisme, beskylder ham for at være i den samme
% `duæ' (p. 27, first line): this fount's italic æ has a closed o-like bowl
% and no separate a-stem, exactly as the certain æ in \emph{Fædrelandet}s, p. 10;
% the bitmap is unique on the page (no JBIG2 substitution).  Both witnesses read `duæ'.  (Batch agent had read `duœ'; overturned at verification, 2026-09-22.)
Grundvildfarelses Vold, der i Middelalderen bragte Menneske%
% --- p. 27 ---
\opage{27}heden til at fastsætte \emph{duæ veritates}. Ti hvem har vel to Sand\-%
heder, om ikke Hr. Wilckens?

Forfatteren slutter sig i denne Lære om Sandhedens Forhold
meget inderligt til Professor R. Nielsen. Naar han undertiden
betegner sin Anskuelse (saaledes f. Eks. den besynderlige Sætning,
at ‹Gud satte Naturlovene til at vidne mod sig selv›) som ‹individuel
Forklaring›, er Meningen heraf maaske den, at han vel ikke i
slige enkelte Punkter, men i hvad han for Resten fremfører, be\-%
tragter sig som Organ for Professor Nielsens Lærdomme. Har han
heri Ret, da finder det Anvendelse paa Danmark nu, hvad Hol\-%
% The colon after `Klim' is italic (600 dpi).
berg siger i \emph{Niels Klim:} ‹Der gives adskillige blandt de euro\-%
pæiske Lærde, der med lige Iver dyrker baade Teologien og
Filosofien; som Filosofer tvivler de om alle Ting, som Teologer
derimod tør de ikke negte nogen Ting›.
\piecerule

% ============================================================
%  Indledning til den rationelle Etik (printed pp. 28--33)
%  Indhold (PDF 9) reads: `Indledning til den rationelle Etik 28'
%  Head as printed: `P. S. V. Heegaard: Indledning til den rationelle Etik'
%    (one line; the initials closely set) / `(1867)'
%  Running head (pp. 29--33) reads: `Indledning til den rationelle Etik'.
%  No drop capital; the text opens with an ordinary indent.
%  p. 33 foot carries the signature line `G. Brandes: Samlede Skrifter.
%  XIII.' and the signature `3' -- not transcribed.
% ============================================================
% --- p. 28 ---
\piecehead{Indledning til den rationelle Etik}{\opage{28}\emph{P. S. V. Heegaard: Indledning til den rationelle Etik}\\[2pt](1867)}
I mere omfattende Betydning af det Ord \emph{Indledning} kan
enhver Forskole, gennem hvilken en Forfatter fører sine
Læsere ind i en Videnskab, kaldes en Indledning til denne.
Baade en Etikens og en Sædelighedens Historie kunde f. Eks.
afgive en Forberedelse til Studiet af Videnskaben om det Gode.
Men af saadanne Propædeutiker gives der for enhversomhelst
Lære kun én uundværlig, nemlig den, der har selve Læren til
Genstand. Den ikke vilkaarligt valgte, men nødvendigt krævede
Indledning til Etiken besvarer det Spørgsmaal: Hvad er Etik?
Etikens eget Problem er: Hvad er det Gode? Da Forfatteren
straks efterat have fastsat Etikens Begreb gaar over til at ud\-%
vikle det Godes Væsen, saa har han ikke givet nogen Indledning
i snevrere og egenlig Forstand.

Men, kunde man sige, lige meget med Navnet: hvad yder
Bogen? Hertil er Svaret, at en Etikens Historie har Dr. Hee\-%
gaard, der intet Hensyn tager til Tidsfølgen, ikke villet give;
han har kun villet benytte Enkeltheder af Historien til at
anskueliggøre de forskellige mulige etiske Principer. At aflede
disse er Opgaven for hans Værk. Da det Inddelingsprincip,
som fastsættes, nemlig Modsætningen imellem det Objektive og
det Subjektive, imidlertid er en Modsætning, der viser sig i
enhver og ikke blot i den etiske Idé, bliver det Gode ikke be\-%
tragtet efter sit ejendommelige Væsen, men kun fra sin almen\-%
logiske Side; de særlig moralske Begreber som Pligt, Dyd, Sam\-%
vittighed tjener efter Planen kun til at oplyse den nævnte Mod\-%
sætning, deres etiske Ejendommelighed kommer da ikke frem, og
% --- p. 29 ---
\opage{29}ligesom Grundmodsætningen, saaledes er ogsaa den imellem
Virkelighed og Væsenlighed og den imellem Umiddelbarhed og
Refleksion, hvorpaa Underinddelingen beror, rent metafysisk.

Men, kunde man sige, her er jo dog et Inddelingsprincip;
er \emph{det} blot fastholdt, da bør man ikke videre gaa i Rette med
Forfatteren. Desværre slippes det lige saa hurtig, som det er
opstillet. Væsenhedsetiken er ikke underinddelt, den subjektive
Etik er slet ikke inddelt, og dog var der aldeles intet i Vejen
for, at de valgte Synspunkter kunde anlægges. Selv om der
% The first dash here is printed as two short rules with a gap (a split
% or double dash, 600 dpi); the second is an ordinary em-dash.  Both
% typed as ---.
ikke gaves historiske Repræsentanter for alle Former --- hvilke
dog let kan findes --- saa krævede Principerkendelsen, paa
hvilken Forfatteren lægger saa megen Vægt, at de i det mindste
alle opstilledes.

Dog, Dr. Heegaard ønsker maaske af Læseren den Ind\-%
rømmelse, at det er mindre vigtigt, om Principet for Under\-%
inddelingen skifter, hvis kun Hovedinddelingsprincipet er fast\-%
holdt. Vi retter da vor Opmærksomhed paa dette. Men det er
vel klart, at skal den subjektivt-objektive Etik være tredje
Led i Rækken, saa kan den ikke have Former under sig, der
paany erklæres for ensidig subjektive som Kierkegaards Etik.
At den ikke destomindre har det, er kun forstaaeligt derved, at
Forfatteren her aldeles har skiftet Inddelingsprincip. De histo\-%
riske Foreteelser, der fremstilles i tredje Afsnit, findes sammen,
ikke fordi deres Etik har et subjektivt-objektivt Præg, hvad den
ikke gennemgaaende har, men fordi de i Modsætning til de
% `Kristeligt-Religiøse': compound hyphen at the line end (cf.
% `Positivt-Religiøse' below); kept.
foregaaende stiller sig i et afgørende Forhold til det Kristeligt-%
Religiøse.

Dette er ikke heldigt, men lad saa være! Vi vil uden
Stridighed følge Forfatteren ogsaa her og blot prøve, hvilke
Former der maa fremkomme, naar saavel det Rationelle som
det Positivt-Religiøse er givne som Forudsætninger. Man kan
da --- for at bruge Forfatterens Udtryk: Tro og Viden ---
enten forkaste Troen for Videns Skyld (som Feuerbach) eller
omvendt Viden til Fordel for Troen (som Kierkegaard), eller
man kan hævde Gyldigheden af begge, og da enten i Enhed
(som Teologien) eller i ubetinget Sondring (som Nielsen). Af disse
Former overspringer Dr. Heegaard aldeles den teologiske Etik:
en Etiker som Schleiermacher kommer ikke med, medens Nielsen,
der slet ingen Etik har skrevet, behandles med megen Vidt\-%
løftighed. Vil Forfatteren nu sige, at den teologiske Etik er for
% --- p. 30 ---
\opage{30}ubetydelig til at tages med, da vil dette ikke være nogen Aar\-%
sag i hans Mund, eftersom han som Repræsentant for en af
Formerne har opstillet Feuerbach, om hvem han selv bemærker,
at denne ikke har noget etisk Princip. Vil han sige, at den
teologiske Etik er en Form, som intet fornuftigt Menneske
behøver at tage Hensyn til, da maa han huske, at det Spørgs\-%
maal endnu er uafgjort, om ikke det samme kan siges om den
Retning, der vil hævde Tro og Viden i absolut Uensartethed.
I ethvert Fald kan man ikke forkaste den teologiske Etik uden
gennem en Kritik; at lukke Døren for den er en Udvej, men
ingen videnskabelig.

Det nye Inddelingsprincip var Modsætningen mellem det
Rationelle og det Positivt-Religiøse, hvilket sidste Dr. Heegaard
haardnakket behandler som enstydigt med det Religiøse i Alminde\-%
lighed, ret som gaves der kun én positiv Religion i én eneste
Bekendelse. Men hvorfra kommer det Religiøse, der falder ned
som en Bombe i Bogens Slutning? Er det nok til en Princip\-%
erkendelse, at vi ikke tager fejl af Subjektivt og Objektivt, Religiøst
og Rationelt, uden at disse Begreber i mindste Maade ses som
Principbestemmelser, som Ytringer af samme Væsen? Forfatteren
% PRINTER'S DEFECT: no stop after `Aandsmagt' (600 dpi: a wide space
% and nothing in it; witness 2's comma is its own); transcribed as printed.
siger blot, at den aabenbarede Religion vitterlig er en Aands\-%
magt Men hvad beviser det? For Principerkendelsen neppe
Meget.

Og hvad beviser det, at Forfatteren kalder det Religiøse
en tvetydig Bestemmelse, eftersom det baade kan tages i rationel
og antirationel Forstand; thi hvorfra kommer det Anti\-%
rationelle? Ja, det kommer fra Aabenbaringen. Og Aaben\-%
baringen? Ja, den er der nu engang. Men er den derfor anti\-%
% Broken type: `dog' has a damaged o (reads `dcg'); transcribed as
% the word.
rationel? Da er der dog saa megen Tanke og Fornuft i det
gamle og nye Testamente, at det i det mindste holder det Mira\-%
kuløse Stangen. Vil man drage det Religiøse ind i Udvik\-%
lingen som bestemt positiv Religionsform, saa maa man tage
Følgerne deraf og med det Samme drage det ind i Under\-%
søgelsen. Vil man ikke dette, men gaar paa ganske umid\-%
delbar Maade ud derfra, saa lad det ske til Gavns, og lad os
faa en Katekismus! Den er ulige nemmere end disse indviklede
Fremstillinger, der dog kun kan paastaa, ikke godtgøre. Men
her er den Dualisme brudt frem, der nu trænger stærkt ind i
vor filosofiske Tilstand, og som viser, at denne hverken er mere
eller mindre end en Krise. Den svækker det Religiøse, der som
% --- p. 31 ---
\opage{31}ubetinget fornuftstridigt bliver en afmægtig Aandskraft, og den svæk\-%
ker Fornuftbuddene, hvilke man berøver de dybeste Elementer,
for at lægge dem over i det Antirationelle, hvad der ikke betyder
meget Andet end at skrinlægge dem. En saakaldt rationel Etik, der
ikke kan og ikke tør paa sin Maade behandle de etiske Spørgs\-%
maal, som enhver Aabenbaring behandler paa sin Maade, er
uden Avlekraft og betydningsløs; den bliver højest en Etik, der
‹lader sig anvende i det borgerlige Samfund og Staten, som de
i Virkeligheden er›, men ikke af et Menneske, i hvem det Etiske
under hvilkensomhelst Form rører sig kraftigt.

Det ældre Inddelingsprincip bringes af Forfatteren i For\-%
bindelse med det ny derved, at den subjektivt-objektive Etik
betegnes som den, der er sig bevidst, at den er Teori og ikke
Praksis, medens Repræsentanterne for den blot objektive og
blot subjektive Etik siges hildede i den Illusion, at Teori og
Praksis uden videre falder sammen --- et stærkt Udtryk i Sand\-%
hed, naar man erindrer, at blandt hine saaledes Fejlende findes
Tænkere som Kant og Fichte. Men til denne ny Modsætning
af Teori og Praksis, dette tredje Inddelingsprincip, der ligesom
det andet kommer himmelfaldent i Slutningen af Bogen, til dette
er der jo ikke taget mindste Hensyn i det tidligere. Man skal
da altsaa tænke de tidligere Former om igen under dette Syns\-%
punkt og saaledes gøre det Arbejde, Forfatteren skulde have
% A speck over the a of `man' (600 dpi: a dot, not an accent); not
% transcribed.
gjort. Han siger selv, at man fra dette Synspunkt har større
Sympati for Aristoteles, Stoikerne og Epikuræerne end for Hegel
og Fichte, men hvorfor blot Sympati? Den hele Ordning bliver
en anden, Hegel og Aristoteles kommer slet ikke mere i Klasse
sammen. Om Forholdet imellem de to sidst fremkomne Ind\-%
delingsprinciper udtaler Dr. Heegaard sig nu desuden aldeles ikke
med Klarhed; det ses vel, at den praktiske Etik kan være baade
rationel og religiøs; men for den første forekommer der ingen
selvstændig Repræsentant, den findes kun som Udviklingstrin
under R. Nielsen.

% PRINTER'S ERROR: `etisko' for `etiske' (1200 dpi: a closed o, no
% crossbar); transcribed as printed.  Witness 2 reads `etiske'.
Forfatteren har altsaa egenlig villet betragte de etisko For\-%
mer under et tredobbelt Synspunkt, nemlig dels efter Modsæt\-%
ningen: subjektiv-objektiv, dels efter den: rationel-religiøs, dels
% `teoretisk-praktisk': the hyphen is printed short and raised, like a
% point (600 dpi); typed as a hyphen.
efter den: teoretisk-praktisk. Hermed er der imidlertid opstillet
tre Inddelingsprinciper. Var nu end Leddelingen nok saa
uklar, saa kunde dog i Kraft af Dialektiken disse tre meget vel
være forenede. Under hver af dem maatte alle Formerne kunne
% --- p. 32 ---
\opage{32}stilles; de maatte gennemtrænge hverandre paa ethvert Punkt.
% The line-end hyphen in `Eksisten-tielle' is printed as a raised
% point (damaged hyphen); a line-break hyphen.
Det burde været vist, hvorledes det Subjektive er det Objektive,
at begge er det Religiøse, begge det Eksistentielle, det Eksisten\-%
tielle det Teoretiske osv., indtil det hele Kresløb var sat i
Virksomhed. Saa kunde Forfatteren med frelst Samvittighed have
gjort sine Sondringer. Nu stritter Formerne uden indre Sammen\-%
hæng, uden al Dialektik udvortes ud fra hinanden, og dog er
Forfatterens Standpunkt det dialektiske, saa dialektisk, at han
endog, hvad Kierkegaard vilde undres ved, lader selve Grænsen
mellem det Rationelle og Antirationelle være af dialektisk Natur.

Et Spørgsmaal paatrænger sig til Slutning, et meget vigtigt,
nemlig dette: Har Forfatteren selv noget Standpunkt? Fordi han
f. Eks. fremstiller Feuerbach som en rationel-ateistisk Etiker,
der ingen Etik har, deraf følger naturligvis ikke, at Forfatteren
er en saadan. Og fordi R. Nielsen fremstilles som den, der
giver en religiøs og rationel Etik, kan man dog ikke forud\-%
sætte som Forfatterens Antagelse, at en saadan har Gyldighed.

Hvad der skulde vise dette, mangler, nemlig en Kritik af
Nielsens Anskuelse. Den mærkelige Kendsgerning finder nemlig
Sted, at da Forfatteren kommer til de Punkter, der som de
afgørende i højeste Grad krævede en Kritik, saa lukker han
meget beskeden sin kritiske Kniv, den han i det Foregaaende
saa kraftigt og tappert har brugt. Professor Nielsen er uden
Sammenligning den betydeligste Repræsentant i den hele sidste
Klasse, og Forfatteren lader ham spille en forunderlig Rolle.
Han er Repræsentant for en særegen Form, men som saadan er
han Repræsentant for Repræsentanter for andre Former, der
% `f. Ex.' here, against `f. Eks.' elsewhere: as printed.
ingen særegen Stemme har faaet, f. Ex. for den rationel-religiøse.
Fremdeles er han en Kritik over det Hele og altsaa ogsaa over
sig selv, og endelig er han Hovedrepræsentant for Etiken, uagtet
han ingen Etik har skrevet.

Hvad vilde man sige, hvis en Kunstdommer holdt Fore\-%
drag om en figurrig Gruppe af Billedhuggerkunst, og efterat
have givet Bedømmelsen af de mindre betydelige Figurer,
forsvandt, naar han skulde bedømme Hovedfiguren, der kastede
det egenlige Lys over det Hele? Kunde han undskylde sig med,
at Hovedfiguren maatte tale for sig selv, og at han ikke vilde
sige Andet, end hvad Hovedfiguren tydelig sagde? Nej, saa
kunde han spare sig hele Foredraget; ti skal hans Tilhørere
kunne bedømme Hovedfiguren, saa kan de ogsaa bedømme
% --- p. 33 ---
\opage{33}hele Gruppen. Enten holder man et Foredrag, og saa holder
man det til Enden, eller man overlader det til Andre, selv at
klare sig Sagen, og saa bliver der slet intet Foredrag af.

Er det altsaa Heegaards Mening, at Nielsens Standpunkt
er hans, saa maa han sige det; det maa ikke være et For\-%
hold, hvorom man kun véd Noget fra en mere eller mindre
paalidelig Haand. Hvilken Anskuelse der ligger til Grund for
Kritiken, er nu en Hemmelighed. Kun dette staar fast, at har
Forfatteren intet andet Standpunkt end Professor Nielsens, saa
er det ligesaa forunderligt at se R. Nielsens Navn mellem de
øvrige Tænkeres, som det vilde være at se Forfatterens eget;
er han derimod ikke Et med Professor Nielsen, da har han
neppe vist stor Selvstændighed eller Uafhængighed ved at lade
dennes Anskuelse staa ubedømt.
\piecerule

% ============================================================
%  Gensvar til Dr. Heegaard (printed pp. 34--39)
%  Indhold (PDF 9) reads: `Gensvar til Dr. Heegaard 34'
%  Head as printed: `Gensvar til Hr. Dr. P. S. V. Heegaard' / `(1867)'
%    -- one line, italic; year roman.
%  Running head (pp. 35--39) reads: `Gensvar til Hr. Dr. P. S. V. Heegaard'
%    (the same as the head).
%  No drop capital; the text opens with an ordinary indent.  One footnote
%  (p. 36).  p. 39 carries only two short paragraph-ends and the closing
%  rule; no dateline or signature.
% ============================================================
% --- p. 34 ---
\piecehead{Gensvar til Dr. Heegaard}{\opage{34}\emph{Gensvar til Hr. Dr. P. S. V. Heegaard}\\[2pt](1867)}
I en Antikritik, hvis Tone ingenlunde er retfærdiggjort ved
Tonen i den Anmeldelse, imod hvilken den er rettet, har Hr.
Dr. P. S. V. Heegaard ikke forsmaaet at benytte det gamle Fif,
hvormed den, der i en aandelig Strid af gode Grunde ikke
ønsker Svar, undertiden forsøger paa at kneble Modstanderen,
og hvis Recept er denne: ‹Dersom Hr. N. N. nu, hvad der
kunde ligne ham, skulde føle en Trang til at have Slutnings\-%
repliken, saa staar det ham frit for etc.›, hvilket er udlagt:
‹Da jeg for min Del intet andet Værge har end at fremsætte
Benegtelser og opkaste Beskyldninger, som ikke kan staa et
Øjeblik længere, end til Gensvar kommer, saa føler jeg en in\-%
derlig Trang til om muligt at afskære min Modstander Ordet.›

Gør det mig nu end ondt ikke at kunne tilfredsstille en saa
dybt følt Trang, saa beder jeg dog Hr. Doktoren være overbe\-%
vist om, at jeg tilfulde forstaar og skatter den. Det er ganske
forklarligt, at Dr. Heegaard, fejret og hyldet af saa sagkyndige
Bedømmere som dem, han har fundet i \emph{Illustreret Tidende} og
\emph{Berlingske}, har begyndt selv at anse sit Arbejde for godt og
grundigt, og at det har været ham en ubehagelig Overraskelse
at finde en Kritiker, der var af modsat Mening; men han burde
have taget denne Kendsgerning med mere Ro; han burde have
stræbt at fatte Tankegangen i min Artikel, i det mindste hvor
den ligger lige for Næsen, og han vilde da ikke have stanset
midt i Læsningen af en Sætningsforbindelse for at jamre over
den første Halvdels Meningsløshed.

% --- p. 35 ---
\opage{35}Her et Eksempel paa, hvorledes Hr. Doktoren har læst. I
min Anmeldelse fandtes følgende Ord: ‹Et Spørgsmaal paatræn\-%
ger sig: Har Forfatteren selv noget Standpunkt? Fordi han
f. Eks. fremstiller Feuerbach som en rationel-ateistisk Etiker,
der ingen Etik har, deraf følger naturligvis ikke, at For\-%
fatteren er en saadan. Og fordi R. Nielsen fremstilles som den,
der giver en religiøs og rationel Etik, kan man dog ikke forud\-%
sætte som Forfatterens Antagelse, at en saadan har Gyldighed.›

Dr. Heegaard river nu den ene Sætning om Feuerbach ud af
dens Sammenhæng, kan ikke fatte Meningen med denne Udta\-%
lelse, eftersom intet fornuftigt Menneske kunde falde paa at
slutte saaledes, og raaber: ‹Hvortil altsaa Bemærkningen?› Hvis
min ærede Modstander nu virkelig ikke i rolig Tilstand for\-%
maaede at indse, hvortil denne Bemærkning tjener, saa vilde
det være forgæves, om jeg forsøgte paa at forklare ham Sagen;
og det er kun, fordi han øjensynlig under Nedskrivningen af
sit Indlæg har befundet sig i en Tilstand af uklar Ilterhed, at
jeg giver ham den nødvendige Oplysning ind med Skeer: De to
opstillede Tilfælde (det med Feuerbach og det med Nielsen) er
ensartede; hvis Intet kan sluttes af det første, saa kan heller
Intet sluttes af det andet, og det er som Følge heraf i selve
Bogen uoplyst, om Forfatteren er \emph{med} Professor Nielsen eller
\emph{imod} Professor Nielsen, eller om han slet ingen Mening har.

Hr. Doktoren viste sig i sin Bog som fuldendt Diplomat,
vogtede sig omhyggeligt for at antyde en selvstændig Opfattelse,
og udførte den Æggedans, ikke i noget Punkt at fjerne sig fra
Professor Nielsen uden dog med en Stavelse at antyde en Over\-%
ensstemmelse med ham, kort sagt: han optraadte som en forsig\-%
tig Mand.

Nu er vel Forsigtighed en uvurderlig Dyd, en Borgmester\-%
dyd, som man siger --- og hvem véd, om ikke Dr. Heegaard
har forfejlet sine Dyders Kald ved endelig at ville være
Filosof! --- men der forekommer dog Tilfælde her i Livet, hvor
man maa opofre den. Der udkræves kun, at En viser sig
saa paatrængende at spørge, om Hr. Doktoren, der ‹er med› i
Andres Udvikling, da slet ingen egen Udvikling har, saa nødes
han til at tilstaa, at han endnu aldrig har tænkt paa Andet end
det Gamle, at sværge paa Lærerens Ord, ja, han anstiller sig
næsten fornærmet over den Mistanke, at han selv kunde have en
Anskuelse, og udleder sin Modtagsomhed overfor Professor Nielsen
% --- p. 36 ---
% Footnote *) (p. 36), set small below a short rule; the whole note is on
% p. 36.  Its type is lightly inked: `Hønen' shows no stroke through the o
% and several of its commas print as points (`mene,' `med,' `Ja,'
% `kontrollere,' `Foredrag,'); read as commas with both witnesses.
\opage{36}‹simpelthen› deraf, at man ikke kritiserer det Standpunkt, hvor\-%
paa man selv befinder sig\fnmark{*)}{Dr. Heegaard taler om ‹at være med i Udviklingen›, omtrent som
Katten og Hønen hos Andersen om at spinde og lægge Æg. Han synes
at mene, at man som \emph{studiosus perpetuus} maa have hørt alle Professor
Nielsens Forelæsninger, dem for Russerne og dem for de Dannede med,
hvis man vil have Lov til at tale med om filosofiske Genstande herhjemme.
Ja, han opfordrer mig endog --- ikke uvittigt --- til at kontrollere, om han
rigtigt har gengivet Professor Nielsens Tankegang efter Foredrag, der endnu
hverken er holdte eller trykte.}. Men hvorfra i al Verden skulde jeg
vel vide Noget om Dr. Heegaards Standpunkt, da den eneste
Vej dertil gik gennem en Slutning, om hvilken Hr. Doktoren selv
(ubesindigt) erklærer, at intet fornuftigt Menneske kunde falde
paa at gøre den?

Nej, Hr. Doktorens Avisartikel yder en Oplysning, om hvil\-%
ken hans Bog ikke giver mindste Vink, og Dr. Heegaard har ved
at meddele den ombyttet sit Overmod med en Beskedenhed uden
Lige. Det være nu langt fra mig at beskylde min Modstander
for ‹Effektjageri›; men mon ikke ligefuldt den Beskedenhed, han
her har lagt for Dagen, skulde gøre sin Virkning? Mon ikke
mangen Fader eller Formynder skulde sige til sin Myndling:
‹O, min Søn! spejl Dig i denne modne Mand! Skønt ikke mere
ganske ung, skønt Doktor i Filosofien og Privatdocent, skønt
Forfatter af en Bog, der indgyder Ærefrygt alene ved sin Tyk\-%
kelse, er han dog saa beskeden endnu ingen anden Mening at
have end den, hans Lærer har.› O! visselig vil Hr. Doktoren,
uden at have tilsigtet det, opnaa en saadan Virkning, og hans
Beskedenhed vil forplante sig fra Slægtled til Slægtled.

Nu er vel Beskedenhed en uskatterlig Dyd, efter Nogles
Mening endog en Ynglings skønneste, men der gives ikke des\-%
mindre Tilfælde her i Livet, hvor en Beskedenhedsytring bliver
Et med en Falliterklæring. Hr. Doktoren er i en stor og be\-%
klagelig Vildfarelse, naar han mener, at man vel behøver et
eget Standpunkt for at skrive en Etik, men ingenlunde for at
skrive en Indledning til samme, og i en anden ligesaa stor, naar
han antager, at min Anke over hans Mangel paa Standpunkt har
sin Grund i en Forveksling af Indledningen til Etiken med Eti\-%
ken selv.

Dr. Heegaard er kæk nok til at forkaste min Bestemmelse
% --- p. 37 ---
% `indrømme. at' (line 17): a point, not a comma, is printed after
% `indrømme' (1200 dpi: a round point, no tail, unlike the comma after
% `Videnskab' below it; witness 1 agrees, witness 2 reads a comma).
% Printer's error or a comma whose tail failed to print; as printed.
\opage{37}af, hvad en Indledning til Etiken vil sige, med et simpelt:
Dette er urigtigt. Vi skal se, med hvad Ret. Det er allerede
besynderligt, at min Modstander vil pukke paa en bestemt For\-%
klaring af Titlen paa hans Bog, ret som var ‹Indledning til Eti\-%
ken› et blandt Videnskabsmænd gængse og antaget Navn paa
en Videnskab. Men det er endnu mere besynderligt, at saa for\-%
sigtig en Forfatter ikke er forsigtigere med den Begrebsbestem\-%
melse, han giver. Naar han siger, at man ikke kan udvikle,
hvad Etik er, uden at forklare, hvad det Gode er, saa er dette
en Selvfølge, absolut forstaaet; men relativt taget er det galt, ti
gøres der ingen Forskel paa Videnskaben om en Videnskab og
paa denne selv, bliver Videnskabernes System til en jævn og
ensartet Taage. Og naar han mener, at Indledningen til Eti\-%
ken bør bestaa i Etikens Metafysik, saa løber han sig fast i
følgende Tvangsvalg: Enten faar han en almindelig Metafysik, der
falder udenfor de særskilte Videnskabers, hvilket er Nonsens,
eller han nødes til at indrømme. at det særegne Fags Metafysik
ikke er nogen selvstændig Videnskab, men maa bestaa af Laane\-%
sætninger fra den almindelige Metafysik. Saaledes forholder
det sig virkelig; derfor er f. Eks. det Skønnes Metafysik heller
aldrig bleven opstillet som en særegen Videnskab. Det Mor\-%
somste er imidlertid, at Dr. Heegaard i Kraft af sin Synsmaade
følgerigtigt maa fordre som filosofisk Propædeutik selve Meta\-%
fysiken, hvorved han dels kommer i Strid med den sunde For\-%
nuft, hvad der her er af mindre Betydning, dels i Strid med
sit eget Orakel, Professor Nielsen, der i Overensstemmelse
med min Opfattelse bestemmer Propædeutiken som den Viden\-%
skab, der besvarer Spørgsmaalet: ‹Hvad er Filosofi?›

Dog jeg skænker Dr. Heegaard hans dristige Paastande om
dette Punkt; lad ham definere Indledningen, som han lyster,
saa Meget staar dog urokket, at hvad en Indledning end er,
maa den være en organisk Bestanddel af det System, hvortil
den fører, og hvoraf den omfattes, og da maa ogsaa samme
Liv, samme Aand, samme Princip pulsere i Indledningen som i
Systemet, og Principløshed være ligesaa forkastelig i hin som i
dette. Hvor kan da en saa beskeden Mand som Dr. Heegaard
tro, at han med sin beskedne Mangel paa Standpunkt er skik\-%
ket til at skrive Indledningen?

Eller har Dr. Heegaard maaske alligevel et Princip, et
% --- p. 38 ---
\opage{38}Princip uden Standpunkt? Hvad har han vel stillet mod
min bestemte Paavisning af hans holdningsløse Vaklen mellem
tre forskellige Inddelingsprinciper? Ikke et Ord, ikke et Bog\-%
% `Saa frejdig kan ellers er': printer's error, `kan' for `han' (read
% at 600 dpi; the k is clear).  Transcribed as printed.
stav! Saa frejdig kan ellers er til at benegte, her overhører
han Kritiken og gør klogt deri. Og er hans Tale om Under\-%
inddelingen maaske heldigere end hans Tavshed om Hoved\-%
punktet? Kan noget mere Bagvendt tænkes end at Hr. Dok\-%
toren opfordrer \emph{mig} til at vise, hvad tvingende Grund der er
til at underinddele Væsenhedsetiken ɔ: til at sønderlemme Fæ\-%
nomenet over det Hjul, hvorpaa han selv har lagt det?

Jeg svarer: Hvis der i Virkelighedens Begreb ligger en tvin\-%
gende Grund til Underinddeling i ‹umiddelbar› og ‹reflekteret›
--- og det maa vel Hr. Doktoren antage, siden han inddeler
saaledes --- da maa der ogsaa ligge en saadan Grund i Væsen\-%
heden. I modsat Tilfælde vilde jo nemlig hine to Bestem\-%
melser falde udenfor denne --- en skøn Ontologi. Efter\-%
trykket hviler paa det Principielle, paa Helheden af alle de
brugte principielle Bestemmelser, der maa udgøre en Enhed.
Dersom det Historiske ikke svarer, saa er det værst for Hr.
Doktoren. Det er ham, ikke mig, der skal faa Stoffet og
Principerne til væsenlig at falde sammen; det er ham, der har re\-%
gistreret, idet han har været ude af Stand til at vise Over\-%
gangen mellem det Historiske og det Principielle, og det vil
neppe lykkes ham at opveje sin Undervægt med den troskyldige
Erklæring, at han ‹ikke har givet saa Lidt endda af Prin\-%
cipudvikling›. Det bliver ikke vissere, fordi han selv for\-%
sikrer det.

Dog, Dr. Heegaard har jo erklæret, at en Kritik af hans
Arbejde, der som min ikke tager alle Spørgsmaalene op for\-%
fra og ikke udfører det uoverkommelige Arbejde at sætte det
Rigtige i det formentlig Urigtiges Sted --- den Kritik ‹er ikke værd
at agte paa›, og denne Sætning har maaske gjort Indtryk paa
En og Anden. Men hvor kan en Forfatter, der efter egen
Tilstaaelse kun har en Helhedsanskuelse paa anden Haand, af
en Dagbladsanmeldelse paa halvtredje Spalte kræve saa me\-%
get Mere, end han selv har givet paa halvfemtehundrede Si\-%
der? Anmelderen har alene leveret Paastande, siger han,
uden Bevis. Det er kun halvt sandt, men lad saa være! Mod
disse stiller han andre Paastande, mindst ligesaa ubeviste, til%
% --- p. 39 ---
\opage{39}syneladende rigtignok uden at være sig bevidst, at han ikke
yder Mere.

Men saa staar endnu det Spørgsmaal tilbage, hvilke
Paastande Bogens tænkende Læser nødsages til at indrømme,
hans eller mine, og jeg overlader det trygt til de Kyndige at
dømme os imellem.
\piecerule

% ============================================================
%  Grundideernes Logik (printed pp. 40--42)
%  Indhold (PDF 9) reads: `Grundideernes Logik 40'
%  Head as printed: `R. Nielsen: Grundideernes Logik'
%    -- one line, italic, and NO YEAR: the short rule follows the title
%    directly (confirmed at 300 and 600 dpi); the only head in pp. 1--42
%    without a year.
%  Running head (pp. 41--42) reads: `R. Nielsen: Grundideernes Logik'
%    (the same as the head).
%  No drop capital; ordinary indent.  No footnotes; no dateline.  p. 42
%  is the last page of the group; p. 43 begins the Dualismen reprint.
% ============================================================
% --- p. 40 ---
\piecehead{Grundideernes Logik}{\opage{40}\emph{R. Nielsen: Grundideernes Logik}}
% line 2: `Professor R. Nielsens' and the title inside the guillemets
% (with its comma) are italic; the guillemets themselves roman.
Med et Par Ord at give en Forestilling om Tankegangen i
\emph{Professor R. Nielsens} ‹\emph{Grundideernes Logik, 2den Del}› lader sig
ikke gøre. At levere en Kritik vilde være endnu vanskeligere.
Men det skal blot siges, at det vilde være en Skam for vort
Fædreland, om et Værk som dette ikke skulde finde talrige Kø\-%
bere blandt det store Publikum, som bestaar af ældre eller
yngre Elever af den udmærkede Universitetslærer og Skribent,
hvem det skyldes. Det sidst udkomne Bind er forfattet med en
logisk Virtuositet, en Skarphed og Bestemthed i Udtryksmaaden,
en allestedsnærværende Livlighed og Flugt i Fremstillingen, der
maa vække den stærkeste Beundring, og som gør Læsningen
ligesaa tiltrækkende som lærerig. Selv den for slige Under\-%
søgelser forholdsvis fremmede Læser vil finde store Partier, der
paa en almenfattelig og tillige ny, original og genifuld Maade
behandler Problemer, som den almindelige Ræsonneren idelig
kommer tilbage til og aldrig ved egne Kræfter formaar at løse.
Saadanne er f. Eks. Spørgsmaalene om Sprogenes Oprindelse
og Betydning, om Ordets Nødvendighed som Udtryk for de
guddommelige Tanker, om Guds Forudviden og Alvidenhed, om
Verdens Skabelse, om Forudtilværelsen osv. Skriftet er gennem\-%
vævet med en om umiskendelig Overlegenhed vidnende Kamp
imod Madvig, imod Sibbern, imod Martensen o. Fl. som Re\-%
præsentanter for énsidige eller forældede Synsmaader, og med
Angreb førte fra nye Standpunkter imod Fortidstænkere som
f. Eks. Kant og Hume, der kaster et ejendommeligt Lys over
disses Fortjenester og Vildfarelser.

% --- p. 41 ---
% `‹Paradoxet›' (x) here against `Paradokset', `Paradokse', `Paradoks'
% (k) everywhere else in the piece: as printed.
\opage{41}Eftersom Stedet her er saa lidet egnet til en kritisk Udta\-%
lelse, skal vi kun om et ganske enkelt Punkt tillade os en Be\-%
mærkning. Da \emph{Grundideernes Logik} under Professor Nielsens
Behandling ikke er bleven nogen strengt afgrænset Videnskab,
og da Forfatteren viser en gennemgaaende Bestræbelse for at
indoptage saa Meget som muligt af det givne religiøse Fore\-%
stillingsindhold, opdukker hist og her Antydninger af Forfatte\-%
rens bekendte Anskuelser om Fornuftomraadets Grænser imod det
Religiøse, skønt Værket endnu kun bærer svage Spor af de
uhyrlige Lærdomme, som Dr. Heegaard i sit sidste Skrift har
fremstillet som Professor Nielsen tilhørende. Hjørnestenen for
hele denne Lære er imidlertid den Kierkegaardske Opfattelse af
‹Paradoxet›, og kunde det blot lykkes at vise, til hvilke for\-%
tvivlede, ja afsindige Følgeslutninger Fastholdelsen af dette Stand\-%
punkt nødvendig maa føre, da er der Haab om, at uden videre
Kritik den hele derpaa opførte Bygning vil falde sammen. Men
et Tilknytningspunkt for et saadant Forsøg har Logikens For\-%
fatter selv givet. Idet han vil godtgøre, at den Kant'ske Kate\-%
gori \emph{Tingen i sig} er et rent Taagebillede, fører han Beviset
saaledes: ‹Kant kan ikke undgaa at fælde en Dom om \emph{Tingen
i sig}. Han maa sige om den, at den \emph{er}. Men er det først
bevist, at Forstandskategorien \emph{Væren} maa være bestemmende
for \emph{Tingen i sig}, da gælder det Samme om alle de øvrige
Kategorier.›

Intet kan være klarere og rigtigere end denne Tankegang.
Men hvor er det da muligt, at den geniale Forfatter af \emph{Grund\-%
ideernes Logik} kan mangle Øje for, at netop selvsamme Bevis\-%
førelse lader sig anvende mod Begrebet ‹det Paradokse›, som
han dog har lyst i Kuld og Køn. Absurditetsfilosofen maa vride
sig, saa Meget han vil, han maa dog tilstaa, at kan han ikke
udsige Andet om Paradokset, dette Uigennemtrængelige, Ube\-%
gribelige, saa maa han dog fastsætte om det, at det \emph{er}. Men
Væren er en Fornuftbestemmelse. Og hører nu ‹Paradokset› til
en Kres, hvori Begrebet ‹Væren› har Gyldighed, saa hører
det ogsaa til en Kres, hvori alle Fornuftbestemmelserne har
Gyldighed. Men hvad ligger heri? Intet Ringere end den himmel\-%
raabende Modsigelse, at det Paradokse som Grundord modtager
lutter fornuftbestemte Omsagn. Vil man undgaa den, da gives der
ingen anden Udvej end at forstumme, ɔ: ikke alene at indeslutte
sig i en Isolerthed, hvori man, skønt man ikke umiddelbart kan
% --- p. 42 ---
\opage{42}meddele sig til Andre, dog forstaar sig selv --- men at forstumme
ubetinget, at blive stum som Fisken er stum, ophøre med at
være Menneske. Dette er Tvangsvalget. Fastholder man ‹det Pa\-%
radokse›, da vil man efterhaanden komme til at optømre et
System som Astronomiens ptolemæiske, hvor den første falske
Gisning bestandig kun kunde opretholdes ved andre ligesaa for\-%
tvivlede, paa samme Maade som det fornuftstridige Paradoks alle\-%
rede har gjort et ‹rationelt Paradoks› (!) nødvendigt. Men da
vil ogsaa en skøn Dag det klare copernikanske Fornuftsystem
staa som det sejrende uden al Kamp ene ved den Alt overvin\-%
dende Magt, der ligger i det Sandes hellige Simpelhed.
\piecerule

\end{document}
